\chapter{Теория}
\label{ch:theory}

\section{Язык программирования Python}

Python – это высокоуровневый, интерпретируемый язык программирования общего назначения, созданный в 1991 году Гвидо ван Россумом. Его основная цель заключалась в создании языка, который был бы одновременно простым для освоения и мощным для решения реальных задач. Python стал широко используемым инструментом благодаря своей читаемости, гибкости и богатому набору встроенных функций.

\subsection*{История языка Python}

Python появился в начале 90-х годов, когда Гвидо ван Россум работал над проектом в Центре математики и информатики (CWI) в Нидерландах. Название языка было вдохновлено британским комедийным шоу ``Monty Python’s Flying Circus'', а не змеями, как многие полагают. Ван Россум хотел создать язык, который объединял бы простоту и читаемость. Он позаимствовал множество идей из языка ABC, который также разрабатывался в CWI, но учел его недостатки, такие как ограниченная масштабируемость. Python с самого начала разрабатывался с ориентацией на поддержку множества парадигм программирования: процедурного, объектно-ориентированного и функционального.

\subsection*{Основные принципы работы Python}

Python базируется на философии, которая отражена в документе "The Zen of Python". Его можно просмотреть, введя в консоли команду \texttt{import this}. Вот некоторые ключевые принципы:

\begin{itemize}
    \item Простота: Код на Python должен быть понятным даже для тех, кто не знаком с проектом;
    \item Читаемость: Python стремится к минимизации сложного и избыточного синтаксиса, чтобы код выглядел естественным;
    \item Явное лучше, чем неявное: Предпочтение отдается кодовым конструкциям, которые понятны с первого взгляда;
    \item Краткость и лаконичность: Python позволяет решать сложные задачи с минимальным количеством кода.
\end{itemize}

\subsection*{Интерпретируемость и кроссплатформенность}

Python является интерпретируемым языком, что означает, что программы выполняются построчно. Это снижает время разработки, так как изменения в коде сразу же можно протестировать. Однако интерпретируемость делает Python медленнее по сравнению с компилируемыми языками. Компилируемые языки, такие как C, C++ или Java, требуют предварительного преобразования исходного кода в машинный код. Этот процесс называется компиляцией и выполняется компилятором. Скомпилированный код сохраняется в исполняемом файле (\texttt{.exe}, \texttt{.out}), который можно запускать без участия исходного кода или дополнительного программного обеспечения, такого как интерпретатор.

\subsubsection*{Преимущества интерпретируемости:}
\begin{itemize}
    \item Отсутствие этапа компиляции: Программисту не нужно вручную компилировать код в машинные инструкции перед выполнением. Это ускоряет цикл разработки, так как изменения в коде можно сразу протестировать;
    \item Портативность: Исходный код Python может выполняться на любой платформе, где установлен интерпретатор Python. Например, одна и та же программа может быть запущена на Windows, macOS и Linux без изменений;
    \item Простота отладки: Интерпретируемый характер Python позволяет легко отслеживать ошибки, так как они обнаруживаются при выполнении конкретной строки кода;
    \item Гибкость: Возможность выполнения кода динамически, даже на лету, делает Python подходящим для таких задач, как скрипты автоматизации, обработка данных и прототипирование.
\end{itemize}

\subsubsection*{Недостатки интерпретируемости:}
\begin{itemize}
    \item Скорость выполнения: Так как код выполняется построчно и требует промежуточного этапа интерпретации, Python работает медленнее, чем компилируемые языки, такие как C++ или Go;
    \item Зависимость от интерпретатора: Для запуска программы требуется установленный интерпретатор Python, что может быть неудобным в некоторых ситуациях.
\end{itemize}

Python также полностью кроссплатформенный. Это означает, что один и тот же код может выполняться на Windows, macOS и Linux, если установлен интерпретатор Python.

\subsection*{Динамическая типизация}

В отличие от статически типизированных языков (например, C++), где тип переменной должен быть указан заранее, Python использует динамическую типизацию. Это означает, что тип переменной определяется автоматически на этапе выполнения. Это делает Python удобным для быстрого прототипирования, но может привести к ошибкам, если типы данных используются непоследовательно.

\subsection*{Автоматическое управление памятью}

Python использует сборщик мусора для управления памятью. Это означает, что разработчику не нужно вручную освобождать память для объектов, которые больше не используются. Сборщик мусора автоматически удаляет ненужные объекты, предотвращая утечки памяти. Однако такая автоматизация накладывает некоторые ограничения. Например, разработчик не имеет прямого контроля над моментом удаления объектов, что иногда может быть критично в задачах с высокой производительностью.

\subsection*{Поддержка нескольких парадигм программирования}

Python поддерживает несколько подходов к программированию:
\begin{itemize}
    \item Процедурное программирование основывается на создании функций и работе с переменными;
    \item Объектно-ориентированное программирование (ООП) код организуется вокруг объектов и классов, которые представляют данные и действия над ними;
    \item Функциональное программирование Python позволяет использовать функции как объекты, поддерживает лямбда-функции, замыкания и модули, такие как \texttt{functools}.
\end{itemize}

\subsection*{Богатая экосистема и стандартная библиотека}

Одним из главных преимуществ Python является его экосистема. Стандартная библиотека Python содержит модули для решения большинства распространенных задач:
\begin{itemize}
    \item Модуль \texttt{os}: работа с операционной системой,
    \item Модуль \texttt{math}: математические операции,
    \item Модуль \texttt{datetime}: работа с датой и временем,
    \item Модуль \texttt{json}: работа с JSON-данными.
\end{itemize}

Дополнительно существует огромный набор сторонних библиотек для веб-разработки (Django, Flask), анализа данных (pandas, NumPy), машинного обучения (scikit-learn, TensorFlow) и многого другого.

\subsection*{Преимущества Python}
\begin{itemize}
    \item Читаемость кода: Простота синтаксиса делает Python популярным выбором для крупных проектов, где требуется участие нескольких разработчиков;
    \item Гибкость: Python используется для веб-разработки, анализа данных, научных исследований, создания игр и даже автоматизации офисных задач;
    \item Быстрое прототипирование: Python позволяет разработчикам быстро проверять идеи благодаря простоте языка и его мощным инструментам;
    \item Активное сообщество: Большое количество ресурсов, документации и форумов облегчает обучение и решение проблем.
\end{itemize}

\subsection*{Недостатки Python}
\begin{itemize}
    \item Cкорость выполнения: Python медленнее, чем языки, такие как C++ или Java, из-за интерпретации кода;
    \item Ограничения многопоточности: GIL (Global Interpreter Lock) ограничивает использование нескольких потоков в одном процессе;
    \item Потребление памяти: Python может использовать больше памяти по сравнению с низкоуровневыми языками из-за динамической типизации и управления памятью.
\end{itemize}

\section{Введение в разработку серверной части веб-приложений}

Современные веб-приложения строятся на основе взаимодействия между клиентской и серверной частями. Серверная часть (бэкенд) выполняет ключевые функции, такие как обработка запросов, управление данными, выполнение бизнес-логики и взаимодействие с базами данных.

\subsection*{Основные понятия}

\subsubsection*{Клиент-серверная архитектура}

Клиент-серверная архитектура — это модель взаимодействия между клиентом (пользовательским устройством, например, браузером или мобильным приложением) и сервером (центральным узлом, который обрабатывает запросы клиента).

Клиент представляет собой программное обеспечение или устройство, которое инициирует взаимодействие с сервером для выполнения определённых задач. Основная роль клиента заключается в отправке запросов на сервер, которые могут включать запросы на получение веб-страниц, данных или выполнение операций, а также в интерпретации полученных от сервера ответов. Это позволяет клиенту отображать необходимую информацию пользователю или выполнять заданные действия. Примерами клиентов могут быть веб-браузеры, мобильные приложения, или специализированные программы автоматизации, которые обрабатывают данные без участия человека.

Сервер, в свою очередь, представляет собой мощную программу или устройство, отвечающее за обработку запросов, поступающих от клиентов. Его основная задача заключается в принятии этих запросов, выполнении соответствующих операций, таких как взаимодействие с базой данных или выполнение бизнес-логики, и отправке результатов обратно клиенту. Сервер играет центральную роль в организации взаимодействия, обеспечивая бесперебойный доступ к данным и услугам. Например, сервер может предоставлять веб-страницы, обрабатывать пользовательские данные или выполнять сложные вычисления по запросу клиента.

Взаимодействие между клиентом и сервером формирует основу работы большинства современных приложений и веб-систем. Клиент выступает интерфейсом для пользователя, предоставляя удобный доступ к функциональности системы, тогда как сервер выполняет всю тяжёлую работу за кулисами, обеспечивая надёжность, безопасность и масштабируемость всей системы.

\subsubsection*{HTTP-протокол}

HTTP (Hypertext Transfer Protocol) — это основной протокол передачи данных в веб-приложениях. Он используется для передачи запросов от клиента к серверу и ответов от сервера к клиенту.

\paragraph{Основные особенности HTTP:}
\begin{itemize}
    \item Протокол "запрос-ответ":
    \begin{itemize}
        \item Клиент отправляет запрос, а сервер отвечает;
        \item Каждый запрос и ответ независимы (HTTP является stateless-протоколом);
    \end{itemize}
    \item Методы HTTP:
    \begin{itemize}
        \item GET — запрос данных с сервера;
        \item POST — отправка данных на сервер;
        \item PUT — обновление данных на сервере;
        \item DELETE — удаление данных на сервере.
    \end{itemize}
    \item Статусы ответов:
    \begin{itemize}
        \item 2xx — успешные ответы;
        \item 4xx — ошибки клиента;
        \item 5xx — ошибки сервера.
    \end{itemize}
\end{itemize}

\section{API (Application Programming Interface)}
API — это интерфейс для взаимодействия между различными программными компонентами. В контексте веб-приложений наиболее часто используется REST API.

\subsection*{REST API:}
\begin{enumerate}
    \item Ресурсы: API работает с ресурсами (например, фильмы, пользователи), которые представляются в виде URL;
    \item Простота и универсальность: REST API основан на использовании стандартных HTTP-методов и кодов статуса;
    \item Форматы данных: Чаще всего данные передаются в формате JSON.
\end{enumerate}

\section{Принципы работы бэкенда}
\subsection*{Обработка запросов}
Основная задача бэкенда — принимать запросы от клиента, обрабатывать их и возвращать ответы. Это включает:
\begin{enumerate}
    \item Получение запроса: Сервер принимает HTTP-запрос от клиента;
    \item Аутентификация и авторизация: Проверка прав доступа пользователя к запрашиваемым данным;
    \item Вызов бизнес-логики: Выполнение операций на сервере, например, вычисления, валидация данных или запуск алгоритмов;
    \item Формирование ответа: Создание ответа с результатами обработки (например, данные из базы данных или сообщение об ошибке);
    \item Отправка ответа: Возвращение результата клиенту в виде HTTP-ответа.
\end{enumerate}

\subsection*{Взаимодействие с базой данных}
Бэкенд играет роль посредника между клиентом и базой данных:
\begin{enumerate}
    \item Чтение данных: Получение данных из базы по запросам клиента;
    \item Запись данных: Сохранение данных, отправленных клиентом;
    \item Обновление данных: Изменение существующих записей;
    \item Удаление данных: Удаление ненужной информации.
\end{enumerate}

\subsection*{Управление бизнес-логикой}
Бизнес-логика — это набор правил и операций, которые управляют работой приложения. Она может включать:
\begin{enumerate}
    \item Проверку условий: Например, возможность добавления фильма только администратором;
    \item Валидацию данных: Убедиться, что данные клиента соответствуют заданным требованиям (например, корректный формат даты выхода фильма);
    \item Выполнение алгоритмов: Например, расчет рекомендаций для пользователя;
    \item Обработка ошибок: Например, возвращение понятных сообщений в случае неудачной операции.
\end{enumerate}

\section{Django как веб-фреймворк}
\subsection*{История и основные особенности Django}
\subsubsection*{История Django}

Django — это популярный веб-фреймворк на языке Python, который был создан в 2005 году разработчиками Эдрианом Хольоватти и Саймоном Уиллисон для управления новостным сайтом. Основной целью создания Django было облегчение разработки сложных веб-приложений за счёт автоматизации рутинных задач. Название фреймворка было выбрано в честь джазового музыканта Джанго Рейнхардта.
С тех пор Django стал одним из самых популярных фреймворков в мире благодаря своей простоте, мощи и философии "не изобретай велосипед". Его используют крупные компании, такие как Instagram, Pinterest и Disqus.

\subsubsection*{Принцип "Батарейки включены" (Batteries Included)}

Одной из ключевых особенностей Django является принцип "Батарейки включены" (batteries included). Это означает, что Django предоставляет все необходимые инструменты "из коробки" для создания веб-приложений. Вам не нужно искать сторонние библиотеки для таких задач, как:
\begin{itemize}
    \item Работа с базой данных (ORM для взаимодействия с базами данных);
    \item Создание административной панели для управления данными;
    \item Настройка роутинга и URL-адресов;
    \item Валидация форм;
    \item Обработка аутентификации и авторизации;
    \item Управление сессиями пользователей;
    \item Поддержка многоязычности и локализации.
\end{itemize}

Принцип "батарейки включены" делает Django идеальным выбором для новичков, так как он минимизирует количество дополнительных зависимостей и упрощает настройку.

\subsubsection*{Шаблон проектирования Model-View-Template (MVT)}

Django использует архитектурный шаблон Model-View-Template (MVT), который похож на классический MVC (Model-View-Controller). Разница в том, что Django перенимает часть функционала контроллера и передаёт его шаблонам и представлениям.
\subsection*{Компоненты MVT:}
В архитектуре MVT (Model-View-Template) каждая часть отвечает за свою задачу и взаимодействует с другими для создания веб-приложения. Model (Модель) описывает данные и их структуру, определяя таблицы базы данных, их связи и ограничения. Она использует встроенный ORM Django для взаимодействия с различными базами данных, что упрощает выполнение операций с данными.

View (Представление) обрабатывает запросы, поступающие от клиента, выполняет логику приложения, включая выборку данных из базы, и формирует ответ. Представление отвечает за связь между моделью и шаблоном, обеспечивая передачу данных в подходящем виде.

Template (Шаблон) отвечает за то, как данные будут отображаться пользователю. С помощью встроенного шаблонизатора Django шаблон динамически генерирует HTML-страницы, предоставляя возможность отображения данных в удобном формате.

\section{Основные компоненты Django}
Django включает множество компонентов для упрощения разработки. Рассмотрим ключевые из них.
\subsection*{Models (Модели)}
Модели — это описания данных, которые ваше приложение использует и сохраняет. Они определяют структуру таблиц в базе данных и позволяют легко взаимодействовать с этими данными с помощью Django ORM.

Особенности моделей:
\begin{itemize}
    \item Каждая модель представляет собой класс Python, унаследованный от \texttt{django.db.models.Model};
    \item Поля модели представляют столбцы таблицы базы данных (например, \texttt{CharField}, \texttt{IntegerField});
    \item Django автоматически создаёт SQL-запросы для работы с базой данных.
\end{itemize}

\subsection*{Views (Представления)}
Представления определяют, какую логику следует выполнить при обработке запросов. Они связывают модели с шаблонами и могут возвращать данные в формате HTML, JSON или других.

В Django существуют два типа представлений:
\begin{enumerate}
    \item Функциональные представления: Определяются как функции Python;
    \item Классовые представления: Предоставляют больше гибкости за счёт наследования и переопределения методов.
\end{enumerate}

\subsection*{Templates (Шаблоны)}
Шаблоны используются для генерации HTML-страниц. Они включают статический HTML и динамические элементы, которые подставляются с помощью встроенного шаблонизатора Django.

\subsection*{URL-конфигурация (роутинг)}
Django предоставляет мощную систему роутинга для сопоставления URL-адресов и представлений. Каждый URL определяется в файле \texttt{urls.py}.

\subsection*{Middleware (Промежуточное ПО)}
Middleware — это компоненты, которые обрабатывают запросы и ответы на уровне фреймворка. Они позволяют добавлять функционал, например:
\begin{itemize}
    \item Аутентификацию;
    \item Логирование запросов;
    \item Управление сессиями;
    \item Обнаружение и предотвращение атак.
\end{itemize}

\section{Преимущества и недостатки Django}
\subsection*{Преимущества Django}
\begin{enumerate}
    \item Быстрый старт: Благодаря принципу "батарейки включены", разработка начинается практически сразу;
    \item Безопасность: Django защищает от распространённых угроз, таких как SQL-инъекции, XSS и CSRF;
    \item Масштабируемость: Django подходит для больших проектов, поддерживая горизонтальное и вертикальное масштабирование;
    \item Сообщество: Большое количество документации, библиотек и активное сообщество;
    \item Административная панель: Встроенная панель администратора для управления данными.
\end{enumerate}

\subsection*{Недостатки Django}
\begin{enumerate}
    \item Монолитность: Django подходит для больших проектов, но для небольших задач он может быть избыточным;
    \item Ограниченная гибкость: Строгая структура может быть неудобной для опытных разработчиков, предпочитающих больше контроля;
    \item Сложности масштабирования: Хотя Django хорошо масштабируется, для очень высоких нагрузок может потребоваться настройка.
\end{enumerate}

\section{Работа с базами данных}
База данных (БД) — это организованная структура для хранения, управления и извлечения информации. Она предназначена для систематизации данных, чтобы облегчить их использование и обеспечить доступ к ним. Данные в БД обычно организованы в таблицы.

\subsection*{Использование ORM в Django}
ORM (Object-Relational Mapping) — это техника для работы с базой данных. ORM позволяет разработчику работать с базой данных как с объектами Python, скрывая детали взаимодействия с реляционной БД.

В Django ORM предоставляет:
\begin{itemize}
    \item Модели, которые описывают структуры данных;
    \item Абстракцию для работы с SQL запросами;
    \item Миграции, позволяющие отслеживать изменения в структуре базы данных.
\end{itemize}

\section{Работа с базами данных}
Веб-приложение для работы с данными о фильмах неизбежно связано с базами данных. Они обеспечивают надёжное хранение, обработку и извлечение информации. В данном разделе мы подробно рассмотрим роль баз данных в веб-приложении, использование ORM в Django, выбор PostgreSQL как системы управления базами данных, а также особенности её интеграции.

\section{Роль баз данных в веб-приложении}
Базы данных занимают центральное место в веб-приложениях. Они обеспечивают:
\begin{itemize}
    \item Хранение структурированных данных: в нашем случае это информация о фильмах, режиссёрах, тегах и комментариях;
    \item Целостность данных: использование ограничений, ключей и связей между таблицами гарантирует их согласованность;
    \item Эффективное извлечение данных: базы данных обеспечивают быстрый доступ к большим объёмам информации через оптимизированные запросы;
    \item Поддержку транзакций: это важно для операций, требующих атомарности, например, добавления фильма вместе с его метаданными и тегами.
\end{itemize}

\section{Использование ORM в Django}
Object-Relational Mapping (ORM) является ключевым компонентом Django, позволяющим разработчикам взаимодействовать с базой данных через модели без необходимости писать SQL-запросы вручную. Это упрощает разработку и делает её более продуктивной.

\subsection*{Связь моделей с таблицами базы данных}
Каждая модель Django представляет таблицу в базе данных. Свойства модели соответствуют столбцам таблицы, а методы модели позволяют работать с данными.

Пример модели для представления информации о фильмах:
\begin{verbatim}
from django.db import models

class Movie(models.Model):
    title = models.CharField(max_length=255)
    description = models.TextField()
    release_date = models.DateField()
    director = models.ForeignKey('Director', on_delete=models.CASCADE)
\end{verbatim}

\section{Создание миграций и управление схемой базы данных}
Миграции в Django используются для управления схемой базы данных. Это удобный инструмент для автоматического создания таблиц, изменения их структуры и миграции данных.

Основные команды для работы с миграциями:
\begin{itemize}
    \item \texttt{python manage.py makemigrations} — создание миграции на основе изменений в моделях;
    \item \texttt{python manage.py migrate} — применение миграций к базе данных.
\end{itemize}

\section{Запросы через ORM}
Django ORM позволяет выполнять CRUD-операции простыми методами:
\begin{itemize}
    \item \textbf{Создание:}
    \begin{verbatim}
    movie = Movie.objects.create(
        title="Inception",
        description="Sci-fi thriller",
        release_date="2010-07-16"
    )
    \end{verbatim}
    \item \textbf{Чтение:}
    \begin{verbatim}
    movies = Movie.objects.filter(
        release_date__year=2010
    )
    \end{verbatim}
    \item \textbf{Обновление:}
    \begin{verbatim}
    movie.title = "Interstellar"
    movie.save()
    \end{verbatim}
    \item \textbf{Удаление:}
    \begin{verbatim}
    movie.delete()
    \end{verbatim}
\end{itemize}

\section{PostgreSQL как выбор для проекта}

\subsection*{Основные особенности PostgreSQL}
PostgreSQL — это мощная объектно-реляционная система управления базами данных с открытым исходным кодом. Её главные особенности:
\begin{itemize}
    \item Поддержка сложных типов данных: JSON, массивы и пользовательские типы;
    \item Расширяемость: возможность добавления пользовательских функций, операторов и типов;
    \item Соответствие стандартам SQL: PostgreSQL строго придерживается стандартов SQL;
    \item Поддержка индексов: включая B-деревья, хэши, GIN и другие типы индексов для ускорения запросов.
\end{itemize}

\subsection*{Достоинства и недостатки PostgreSQL}

\textbf{Достоинства:}
\begin{itemize}
    \item Надёжность и стабильность в работе с большими объёмами данных;
    \item Широкий функционал для аналитики, обработки данных и написания сложных запросов;
    \item Поддержка транзакций и механизмов блокировок для обеспечения целостности данных.
\end{itemize}

\textbf{Недостатки:}
\begin{itemize}
    \item Относительно высокая сложность настройки;
    \item Потребление ресурсов на больших нагрузках может превышать аналоги, такие как MySQL.
\end{itemize}

\section{СУБД проекта}
Выбор PostgreSQL для нашего проекта был обусловлен несколькими факторами:
\begin{enumerate}
    \item Поддержка расширенных возможностей: работа с JSON пригодилась для хранения и обработки тегов и комментариев;
    \item Сообщество и документация: доступ к множеству материалов по настройке и оптимизации;
    \item Оптимизация запросов: возможность использования индексов для ускорения фильтрации и поиска;
    \item Интеграция с Django: библиотека Psycopg2 упрощает подключение к PostgreSQL и её использование в Django ORM.
\end{enumerate}

Работа с базами данных является ключевым аспектом любого веб-приложения. Использование Django ORM значительно упрощает процесс взаимодействия с базой данных, а выбор PostgreSQL гарантирует надёжность и высокую производительность системы. Это позволяет создавать гибкие и масштабируемые приложения, которые легко поддерживать и развивать в дальнейшем.

\section{Схема базы данных и нормализация данных}
Базы данных являются основой практически любого современного веб-приложения. Для нашего сервиса, предназначенного для хранения данных о фильмах, схема базы данных играет ключевую роль в организации данных и обеспечении их целостности. В данном разделе рассматриваются логическая структура базы данных, её связи и принципы нормализации, которые мы применили в проекте.

\section{Логическая структура базы данных для хранения информации о фильмах}
Основная задача базы данных нашего проекта — хранение информации о фильмах, жанрах, режиссёрах, тегах и комментариях. Чтобы эффективно структурировать данные и упростить доступ к ним, мы реализовали следующую схему базы данных:
\begin{enumerate}
    \item \textbf{Таблица “Movies” (Фильмы):}
    \begin{itemize}
        \item Поля: 
        \begin{itemize}
            \item \texttt{id} (уникальный идентификатор фильма),
            \item \texttt{title} (название фильма),
            \item \texttt{release\_date} (дата выхода),
            \item \texttt{description} (описание фильма),
            \item \texttt{director\_id} (связь с режиссёром).
        \end{itemize}
        \item Описывает основные характеристики каждого фильма.
    \end{itemize}
    \item \textbf{Таблица “Directors” (Режиссёры):}
    \begin{itemize}
        \item Поля: 
        \begin{itemize}
            \item \texttt{id} (уникальный идентификатор режиссёра),
            \item \texttt{name} (имя режиссёра).
        \end{itemize}
        \item Связь с таблицей “Movies” — один ко многим (один режиссёр может снять множество фильмов).
    \end{itemize}
    \item \textbf{Таблица “Genres” (Жанры):}
    \begin{itemize}
        \item Поля: 
        \begin{itemize}
            \item \texttt{id} (уникальный идентификатор жанра),
            \item \texttt{name} (название жанра).
        \end{itemize}
        \item Связь с таблицей “Movies” через промежуточную таблицу (многие ко многим, так как один фильм может иметь несколько жанров, а один жанр может относиться к нескольким фильмам).
    \end{itemize}
    \item \textbf{Таблица “Tags” (Теги):}
    \begin{itemize}
        \item Поля: 
        \begin{itemize}
            \item \texttt{id} (уникальный идентификатор тега),
            \item \texttt{name} (название тега).
        \end{itemize}
        \item Также связана с таблицей “Movies” через промежуточную таблицу.
    \end{itemize}
    \item \textbf{Таблица “Comments” (Комментарии):}
    \begin{itemize}
        \item Поля: 
        \begin{itemize}
            \item \texttt{id} (уникальный идентификатор комментария),
            \item \texttt{movie\_id} (ссылка на фильм, к которому относится комментарий),
            \item \texttt{user\_id} (ссылка на пользователя, оставившего комментарий),
            \item \texttt{text} (текст комментария),
            \item \texttt{created\_at} (дата и время создания).
        \end{itemize}
        \item Связь с таблицей “Movies” — один ко многим (один фильм может иметь множество комментариев).
    \end{itemize}
    \item \textbf{Таблица “Users” (Пользователи):}
    \begin{itemize}
        \item Поля: 
        \begin{itemize}
            \item \texttt{id} (уникальный идентификатор пользователя),
            \item \texttt{username} (имя пользователя),
            \item \texttt{email} (электронная почта пользователя),
            \item \texttt{password} (хэшированный пароль).
        \end{itemize}
        \item Пользователи могут оставлять комментарии к фильмам.
    \end{itemize}
\end{enumerate}

\section{Пример схемы данных и пояснение её структуры}
В основе структуры базы данных лежат понятия нормализованных таблиц и установленных между ними связей. Пример взаимодействий:
\begin{itemize}
    \item Фильм связан с жанрами через промежуточную таблицу \texttt{MovieGenres}, где хранятся два внешних ключа: \texttt{movie\_id} и \texttt{genre\_id};
    \item Таблица пользователей Users объединяется с таблицей комментариев \texttt{Comments} через внешние ключи, указывающие, какой пользователь оставил какой комментарий;
    \item Таблица режиссёров \texttt{Directors} позволяет избежать дублирования информации о режиссёре при добавлении нескольких фильмов одного автора.
\end{itemize}
Подобная структура обеспечивает гибкость в управлении данными и минимизирует дублирование.

\section{Принципы нормализации базы данных}
При проектировании схемы базы данных мы применили методы нормализации, чтобы избежать избыточности данных и обеспечить их целостность. Рассмотрим основные нормальные формы, использованные в проекте:

\section{Первая нормальная форма (1NF)}
\begin{itemize}
    \item Все данные хранятся в атомарных полях.
    \item В таблицах нет повторяющихся строк, и каждое поле содержит только одно значение.
    \item Например, в таблице ``Movies'' жанры не хранятся в виде списка, а вынесены в отдельную таблицу ``Genres'' с промежуточной таблицей для связи.
\end{itemize}

\section{Вторая нормальная форма (2NF)}
\begin{itemize}
    \item Все неключевые атрибуты зависят от первичного ключа.
    \item В таблице ``Comments'' информация о фильме и пользователе хранится как ссылки на соответствующие таблицы, что позволяет исключить дублирование данных.
\end{itemize}

\section{Третья нормальная форма (3NF)}
\begin{itemize}
    \item Удалены транзитивные зависимости между неключевыми атрибутами.
    \item Например, информация о режиссёре хранится отдельно в таблице ``Directors'', а не дублируется в каждом фильме.
\end{itemize}

\section{Избежание дублирования данных и обеспечение целостности информации}
Одной из ключевых задач нормализации является уменьшение избыточности данных. В проекте этого удалось добиться за счёт:
\begin{itemize}
    \item Использования связей между таблицами (внешние ключи), что исключило повторение одних и тех же данных в разных местах.
    \item Применения промежуточных таблиц для отношений многие ко многим (например, между фильмами и жанрами).
    \item Валидации данных на уровне базы и приложения для предотвращения введения некорректной информации.
\end{itemize}

Схема базы данных для сервиса фильмов была разработана с учётом принципов нормализации, что обеспечило её гибкость, надёжность и простоту управления. Логическая структура и реализованные связи между таблицами не только позволили эффективно обрабатывать запросы пользователей, но и создали основу для масштабирования системы и добавления нового функционала в будущем.

\section{Интеграция фронтенда и бэкенда}
Современные веб-приложения состоят из двух ключевых компонентов: фронтенда, который отвечает за взаимодействие с пользователем, и бэкенда, обеспечивающего управление данными и обработку запросов. Эти два компонента тесно связаны через стандартизированные механизмы обмена информацией, такие как REST API. В данном разделе рассматриваются теоретические и практические аспекты взаимодействия фронтенда и бэкенда в рамках нашего проекта.

\section{Роль фронтенда в веб-приложении}
Фронтенд играет важнейшую роль в предоставлении пользователю удобного интерфейса для работы с приложением. Он отвечает за визуальное представление данных, полученных с бэкенда, а также за сбор и передачу данных, вводимых пользователем. Основная цель фронтенда — сделать взаимодействие пользователя с приложением интуитивно понятным, быстрым и приятным.

В рамках нашего проекта фронтенд предоставляет возможность:
\begin{itemize}
    \item Просматривать информацию о фильмах, режиссёрах, жанрах и тегах.
    \item Добавлять, редактировать и удалять комментарии или метки к фильмам.
    \item Получать персонализированные рекомендации на основе данных, хранящихся на бэкенде.
\end{itemize}

\section{Технологии, используемые для создания фронтенда}
Разработка фронтенда основана на использовании технологий HTML, CSS и JavaScript, которые обеспечивают структуру, стилизацию и логику работы пользовательского интерфейса. В современных веб-приложениях чаще всего применяются также фронтенд-фреймворки, такие как React или Vue.js, которые позволяют:
\begin{itemize}
    \item Создавать динамические и адаптивные интерфейсы с использованием компонентов.
    \item Оптимизировать загрузку данных и обработку событий на стороне клиента.
    \item Упрощать управление состоянием приложения при взаимодействии с API.
\end{itemize}

\section{Организация взаимодействия между фронтендом и бэкендом}
\subsection*{REST API как способ обмена данными}
REST API (Representational State Transfer) является стандартным подходом к созданию интерфейсов для обмена данными между фронтендом и бэкендом. Основные преимущества использования REST API в нашем проекте:
\begin{itemize}
    \item Универсальность: API можно использовать для работы с любым клиентом (веб, мобильные приложения).
    \item Масштабируемость: REST API легко расширяется новыми ресурсами и функционалом.
    \item Стандартизация: использование стандартных методов HTTP (``GET'', ``POST'', ``PUT'', ``DELETE'') упрощает понимание и интеграцию.
\end{itemize}

\section{Работа с форматами JSON для передачи данных}
Фронтенд и бэкенд обмениваются данными в формате JSON (JavaScript Object Notation). JSON был выбран в качестве основного формата из-за его следующих преимуществ:
\begin{itemize}
    \item Читабельность для человека: структура данных интуитивно понятна и проста в использовании.
    \item Совместимость: JSON поддерживается большинством языков программирования и библиотек.
    \item Лёгкость обработки: данные в JSON можно быстро парсить и преобразовывать в объекты JavaScript.
\end{itemize}

\section{Использование запросов (fetch, axios) для обращения к бэкенду}
Для работы с API во фронтенде часто используются JavaScript-библиотеки, такие как ``fetch'' или ``axios''. Эти инструменты позволяют:
\begin{itemize}
    \item Выполнять HTTP-запросы к бэкенду с минимальными усилиями.
    \item Управлять ошибками и временем ожидания запросов.
    \item Гибко конфигурировать заголовки, параметры и тело запросов.
\end{itemize}

Пример запроса данных о фильме с использованием ``axios'':
\begin{verbatim}
axios.get('https://example.com/api/movies/1/')
  .then(response => {
    console.log(response.data);
  })
  .catch(error => {
    console.error('Ошибка при получении данных:', error);
  });
\end{verbatim}

\section{Пример взаимодействия фронтенда и бэкенда}
Рассмотрим сценарий получения данных о фильме и их отображения на странице. Фронтенд отправляет запрос к API с помощью библиотеки ``fetch'' или ``axios''. В ответ бэкенд возвращает JSON с информацией о фильме, которую фронтенд обрабатывает и отображает пользователю:
\begin{enumerate}
    \item Фронтенд отправляет запрос:
    \begin{verbatim}
    fetch('https://example.com/api/movies/1/')
      .then(response => response.json())
      .then(data => {
        document.getElementById('movie-title').innerText = data.title;
        document.getElementById('movie-description').innerText = data.description;
      });
    \end{verbatim}
    \item Бэкенд, реализованный на Django REST Framework, отвечает данными:
    \begin{verbatim}
    {
      "id": 1,
      "title": "Inception",
      "description": "A skilled thief is given a chance at redemption...",
      "director": "Christopher Nolan",
      "tags": ["thriller", "sci-fi"]
    }
    \end{verbatim}
    \item Фронтенд обновляет DOM, отображая название, описание и другие детали фильма.
\end{enumerate}

\section{Использование Postman для тестирования API}
На этапе разработки мы активно использовали Postman — инструмент для тестирования API. Он позволяет:
\begin{itemize}
    \item Создавать и сохранять HTTP-запросы.
    \item Тестировать методы API (``GET'', ``POST'', ``PUT'', ``DELETE'').
    \item Проверять структуру возвращаемых данных.
    \item Выявлять ошибки на стороне бэкенда до подключения фронтенда.
\end{itemize}

Пример использования Postman: запрос данных о фильме с помощью метода ``GET'' по адресу https://example.com/api/movies/1/. В ответ возвращаются те же данные в формате JSON, которые отображаются во фронтенде.

\section{Эффективная интеграция фронтенда и бэкенда}
Эффективная интеграция фронтенда и бэкенда является неотъемлемой частью современного веб-разработки. REST API обеспечивает стандартизированный и масштабируемый способ обмена данными, а использование инструментов, таких как Postman, помогает тестировать и отлаживать взаимодействие. В результате проект становится более надёжным и легко расширяемым для поддержки новых функций.

\section{Создание REST API}
REST (Representational State Transfer) — это архитектурный стиль, который используется для разработки веб-сервисов. REST определяет ряд принципов, направленных на создание простого, масштабируемого и стандартизированного подхода к взаимодействию между клиентом и сервером. В данном разделе описываются ключевые аспекты REST, реализация REST API с использованием Django REST Framework (DRF), а также важность унифицированного подхода в построении взаимодействия между сервисами.

\section{Что такое REST и его ключевые принципы}
REST основывается на следующих ключевых принципах:
\begin{itemize}
    \item Ресурсы: Все данные и функционал сервиса представляются в виде ресурсов, каждый из которых идентифицируется уникальным URI.
    \item CRUD-операции: REST API использует стандартные методы HTTP (‘GET’, ‘POST’, ‘PUT’, ‘DELETE’), которые соответствуют операциям чтения, создания, обновления и удаления данных.
    \item Идемпотентность: Некоторые HTTP-методы (‘PUT’, ‘DELETE’) являются идемпотентными, что означает, что их повторное выполнение приведет к одному и тому же результату.
    \item Состояние клиента: RESTful-сервисы не хранят состояние клиента на сервере, что делает их масштабируемыми и независимыми от состояния сессии.
    \item Единообразие интерфейсов: Взаимодействие клиента и сервера стандартизировано, что упрощает интеграцию разных сервисов.
\end{itemize}

Использование REST API позволяет создать универсальный и стандартизированный способ обмена данными, который может быть использован различными клиентами: веб-приложениями, мобильными устройствами и другими сервисами.

\section{Реализация REST API с Django REST Framework (DRF)}
Django REST Framework предоставляет удобные инструменты для создания REST API. DRF значительно упрощает разработку и поддерживает все ключевые аспекты REST.

\section{Сериализаторы (Serializers) и их роль}
Сериализаторы — это классы, которые используются для преобразования данных между Python-объектами и форматами, такими как JSON или XML. В нашем проекте они играют ключевую роль при обработке запросов и ответов. Основные функции сериализаторов:
\begin{itemize}
    \item Преобразование данных моделей в формат JSON для передачи на фронтенд.
    \item Проверка входных данных, полученных в запросах, и преобразование их в формат, подходящий для работы в Python.
\end{itemize}

\section{DRF предоставляет стандартные классы}
DRF предоставляет стандартные классы, такие как ModelSerializer, который упрощает создание сериализаторов, автоматически генерируя поля и методы на основе модели.

\section{ViewSets и маршрутизация запросов}
ViewSets — это высокоуровневые классы, которые объединяют в себе обработку CRUD-операций для конкретного ресурса. Вместо того чтобы писать отдельные методы для каждой операции (‘GET’, ‘POST’ и т. д.), ViewSets автоматически обрабатывают их в соответствии с настройками.

DRF также предоставляет механизм маршрутизации запросов с использованием routers. Это позволяет легко связывать URL-адреса с соответствующими ViewSets, минимизируя вероятность ошибок. Например:
\begin{verbatim}
from rest_framework.routers import DefaultRouter
from .views import MovieViewSet

router = DefaultRouter()
router.register(r'movies', MovieViewSet)
urlpatterns += router.urls
\end{verbatim}
Такой подход значительно упрощает настройку маршрутов и поддерживает читабельность кода.

\section{Понятие аутентификации и авторизации}
REST API часто используется для работы с конфиденциальными данными, что требует надежных механизмов аутентификации и авторизации. Django REST Framework поддерживает множество способов аутентификации:
\begin{itemize}
    \item Токены: Генерируются для каждого пользователя и передаются в заголовках HTTP-запросов. Библиотека rest\_framework.authtoken предоставляет встроенный механизм токенной аутентификации.
    \item Сессии: Используются для хранения данных аутентификации между запросами в пределах одной сессии.
    \item JWT (JSON Web Tokens): Позволяет создавать защищенные токены с полезной нагрузкой. JWT токены являются компактными и подходят для использования в веб- и мобильных приложениях.
\end{itemize}

Пример настройки токенной аутентификации:
\begin{verbatim}
INSTALLED_APPS += [
    'rest_framework',
    'rest_framework.authtoken',
]

REST_FRAMEWORK = {
    'DEFAULT_AUTHENTICATION_CLASSES': [
        'rest_framework.authentication.TokenAuthentication',
    ],
}
\end{verbatim}

\section{Преимущества использования REST API}
\begin{enumerate}
    \item Универсальность: REST API предоставляет единый способ взаимодействия, который подходит для разных типов клиентов (веб, мобильные приложения, внешние сервисы).
    \item Масштабируемость: Отсутствие сохраненного состояния на сервере упрощает масштабирование приложения.
    \item Удобство разработки: Простая структура запросов и поддержки делает REST API легко понятным для разработчиков.
    \item Проверка запросов: Инструменты, такие как Postman, позволяют тестировать и проверять API перед интеграцией с клиентом.
\end{enumerate}

REST API в нашем проекте позволил фронтенду и бэкенду эффективно взаимодействовать, сохраняя при этом гибкость, расширяемость и удобство поддержки приложения.

\section{Сравнение Django с FastAPI}

\section{Обзор FastAPI}
FastAPI — это современный веб-фреймворк для создания API, разработанный на Python. Он был создан с акцентом на высокую производительность, простоту использования и поддержку современных возможностей Python, таких как аннотации типов. FastAPI быстро завоевал популярность среди разработчиков, благодаря своей скорости и удобству для написания мощных API с минимальными усилиями.

Основные характеристики FastAPI:
\begin{itemize}
    \item Использование аннотаций типов Python: Благодаря встроенной поддержке аннотаций типов FastAPI позволяет автоматически проверять входные данные, генерировать документацию API и ускорять разработку.
    \item Высокая производительность: Построенный на базе ASGI (Asynchronous Server Gateway Interface), FastAPI обеспечивает асинхронную обработку запросов, что делает его чрезвычайно быстрым и подходящим для высоконагруженных систем.
    \item Интеграция со стандартами: Встроенная поддержка OpenAPI и JSON Schema облегчает создание самодокументируемых API.
    \item Асинхронность: Полная поддержка async/await, что позволяет более эффективно работать с большими объемами данных и долгими запросами.
\end{itemize}

\section{Плюсы FastAPI}
\begin{enumerate}
    \item Скорость разработки: Использование аннотаций типов Python позволяет разрабатывать API быстрее благодаря автоматической валидации данных и генерации документации.
    \item Высокая производительность: Благодаря своей асинхронной архитектуре FastAPI быстрее большинства других фреймворков, таких как Django. Он идеально подходит для микросервисов и систем с высоким числом одновременных запросов.
    \item Автоматическая документация: FastAPI генерирует документацию в формате OpenAPI прямо из кода, используя инструменты Swagger UI и ReDoc. Это особенно полезно для крупных проектов, где требуется согласованность API.
    \item Лёгкость освоения: Для разработчиков, знакомых с современными подходами разработки на Python, переход на FastAPI происходит быстро.
    \item Поддержка асинхронного кода: Встроенная поддержка async/await обеспечивает эффективную работу в проектах с большим количеством запросов к внешним системам, например, базам данных или сторонним API.
\end{enumerate}

\section{Минусы FastAPI}
\begin{enumerate}
    \item Недостаток встроенных модулей: В отличие от Django, который предоставляет модули для работы с пользователями, сессиями, авторизацией, ORM и админ-панелью, в FastAPI многое приходится разрабатывать с нуля. Это увеличивает время разработки для проектов, где эти функции необходимы.
    \item Относительно молодая экосистема: FastAPI появился в 2018 году, поэтому его экосистема меньше и менее зрелая, чем у Django. Найти готовые решения для частных задач бывает сложнее.
    \item Сложности для новичков: Несмотря на дружественность к разработчикам, начинающим может быть непросто разобраться в асинхронном программировании, которое активно используется в FastAPI.
    \item Сложность интеграции с некоторыми инструментами: Многие сторонние инструменты пока не имеют официальной поддержки для FastAPI, что иногда требует написания кастомных решений.
\end{enumerate}

\section{Сравнение Django с FastAPI}
Для проекта хранения данных о фильмах Django был выбран вместо FastAPI по нескольким причинам:

\begin{enumerate}
    \item \textbf{Богатая документация}: Django имеет обширную и детально проработанную документацию. Это позволяет легко найти примеры и решения практически для любой задачи. Для учебного проекта с ограниченным временем на освоение нового инструмента это стало значительным преимуществом.
    \item \textbf{Встроенные модули}: Django предлагает множество встроенных модулей, таких как система аутентификации, админ-панель, инструмент работы с формами, встроенный ORM и другие. Это значительно сокращает время разработки, так как многие вещи уже реализованы и требуют минимальных настроек.
    \item \textbf{Простота интеграции типовых решений}: Django «из коробки» обеспечивает интеграцию с базами данных, управление миграциями и создание CRUD-операций. Это позволило сконцентрироваться на изучении бизнес-логики проекта, не отвлекаясь на мелкие технические детали.
    \item \textbf{Распространённость в обучении}: Django часто используется в качестве учебного инструмента из-за своей понятной структуры. Команда проекта уже имела некоторый опыт работы с этим фреймворком, что также поспособствовало выбору.
\end{enumerate}


FastAPI — мощный инструмент для создания асинхронных и высокопроизводительных API. Однако, для учебного проекта, где требуется реализация стандартного функционала с минимальными трудозатратами, Django оказался более подходящим выбором. Он предложил готовые решения для большинства задач, богатую документацию и зрелую экосистему, что позволило команде сосредоточиться на бизнес-логике приложения.

\section{Контейнеризация}

\subsection*{Что такое контейнеризация?}
Контейнеризация — это технология упаковки программного обеспечения вместе со всеми его зависимостями в единый контейнер. Этот контейнер можно запустить на любой системе с установленной поддержкой контейнеров. Контейнеры обеспечивают изоляцию приложения от операционной системы, что делает приложение более стабильным и предсказуемым в различных средах.

\subsection*{Зачем нужна контейнеризация?}
В современном программировании программы становятся всё сложнее, завися от множества библиотек, конфигураций и внешних компонентов. Эти зависимости могут работать по-разному в разных операционных системах или окружениях. Контейнеризация решает эти проблемы, предоставляя стандартизированное и изолированное окружение для запуска приложений.

Основные задачи контейнеризации:
\begin{enumerate}
    \item \textbf{Изоляция окружения}: Разные проекты могут иметь конфликтующие зависимости (например, разные версии Python или библиотек). Контейнеры позволяют избежать таких конфликтов.
    \item \textbf{Упрощение развёртывания}: Контейнер содержит всё необходимое: код приложения, библиотеки, конфигурационные файлы и окружение.
    \item \textbf{Масштабируемость}: Контейнеры можно легко копировать, что упрощает запуск нескольких экземпляров приложения для обработки большого числа запросов.
    \item \textbf{Стабильность}: Контейнеры гарантируют, что приложение будет работать одинаково на разных компьютерах и серверах, независимо от окружения.
\end{enumerate}

\section{Преимущества и недостатки контейнеризации}
Преимущества:
\begin{enumerate}
    \item \textbf{Универсальность}: Контейнеры можно запустить на любой операционной системе, поддерживающей контейнеризацию (Linux, Windows, macOS).
    \item \textbf{Повышение производительности}: Контейнеры легче виртуальных машин, что уменьшает затраты ресурсов и ускоряет их запуск.
    \item \textbf{Масштабируемость}: Контейнеры позволяют легко развернуть несколько экземпляров одного приложения для распределения нагрузки.
    \item \textbf{Быстрое развёртывание}: Контейнеры упрощают установку приложений на разные машины, что особенно полезно для командной работы.
    \item \textbf{Изоляция приложений}: Каждый контейнер изолирован, что предотвращает конфликты между приложениями, работающими в одном окружении.
\end{enumerate}

Недостатки:
\begin{enumerate}
    \item \textbf{Сложность настройки}: Для работы с контейнерами требуется изучение новых инструментов, таких как Docker или Kubernetes.
    \item \textbf{Ограничения изоляции}: Контейнеры используют ядро хост-системы, что может привести к некоторым ограничениям в безопасности.
    \item \textbf{Управление множеством контейнеров}: При большом количестве контейнеров управление ими вручную становится сложным. Здесь требуются системы оркестрации, такие как Kubernetes.
\end{enumerate}

\section{Docker и его роль в контейнеризации}
Docker — это популярная платформа для создания, развертывания и управления контейнерами. Он позволяет разработчикам упаковывать приложение вместе с его зависимостями в контейнер, который можно запустить практически на любой машине.

\section{Основные компоненты Docker}
\begin{itemize}
    \item \textbf{Docker Engine} — ядро Docker, отвечающее за создание и запуск контейнеров.
    \item \textbf{Docker Hub} — репозиторий, где разработчики могут хранить и делиться готовыми контейнерными образами.
    \item \textbf{Dockerfile} — скрипт для автоматизации процесса сборки контейнерного образа.
\end{itemize}

\section{Docker Compose: управление сервисами}
Docker Compose — инструмент для описания и управления многокомпонентными приложениями. Он позволяет разработчикам описывать зависимые сервисы в одном файле docker-compose.yml, что упрощает управление контейнерами.

Пример проекта с использованием Docker Compose:
\begin{enumerate}
    \item \textbf{Backend (Django)}: отвечает за API и бизнес-логику.
    \item \textbf{Database (PostgreSQL)}: база данных для хранения информации.
    \item \textbf{NGINX}: прокси-сервер для маршрутизации запросов и выдачи статических файлов.
\end{enumerate}

\begin{verbatim}
version: "3.9"
services:
  web:
    build: .
    ports:
      - "8000:8000"
    volumes:
      - .:/app
    depends_on:
      - db
  db:
    image: postgres
    environment:
      POSTGRES_USER: user
      POSTGRES_PASSWORD: password
      POSTGRES_DB: mydatabase
  nginx:
    image: nginx
    ports:
      - "80:80"
    volumes:
      - ./nginx.conf:/etc/nginx/nginx.conf
\end{verbatim}

Такой подход упрощает локальную разработку: достаточно одной команды \texttt{docker-compose up}, чтобы запустить все сервисы проекта.

\section{Оркестрация контейнеров}
Когда проект вырастает до масштабов, где необходимо управлять множеством контейнеров, используется оркестрация. Kubernetes — самый популярный инструмент для этой задачи. Оркестрация позволяет:
\begin{itemize}
    \item Автоматизировать развёртывание контейнеров.
    \item Обеспечивать их стабильную работу.
    \item Масштабировать проект.
\end{itemize}

\section{Контейнеризация и командная работа}
При работе в команде Docker существенно ускоряет процесс разработки:
\begin{itemize}
    \item Каждый разработчик может настроить окружение проекта за считанные минуты, выполнив несколько команд.
    \item Единое окружение позволяет избежать ошибок, вызванных различиями в настройках компьютеров участников.
    \item Контейнеризация делает приложение готовым к развертыванию на любых серверах, включая облачные платформы, такие как AWS, Google Cloud или Azure.
\end{itemize}

Контейнеризация позволяет сделать процесс разработки и развёртывания приложений быстрым, предсказуемым и масштабируемым. Docker и Docker Compose стали неотъемлемой частью современных рабочих процессов, обеспечивая изоляцию окружения, лёгкость масштабирования и стабильность работы приложений. В рамках командной работы использование Docker облегчает настройку окружения, сводя вероятность ошибок к минимуму.

\hypertarget{ux432ux432ux435ux434ux435ux43dux438ux435-ux432-ux43cux430ux448ux438ux43dux43dux43eux435-ux43eux431ux443ux447ux435ux43dux438ux435}{%
\section{Введение в машинное
обучение}\label{ux432ux432ux435ux434ux435ux43dux438ux435-ux432-ux43cux430ux448ux438ux43dux43dux43eux435-ux43eux431ux443ux447ux435ux43dux438ux435}}

Машинное обучение --- это наука о нахождении закономерностей в данных с
помощью алгоритмов, позволяющих строить модели на основе примеров. Этот
подход полезен там, где нет точных решений, но есть накопленные данные,
например, для прогнозирования прибыли ресторанов или поведения клиентов.

\hypertarget{ux437ux430ux434ux430ux447ux438-ux43cux430ux448ux438ux43dux43dux43eux433ux43e-ux43eux431ux443ux447ux435ux43dux438ux44f}{%
\subsection*{Задачи машинного
обучения}\label{ux437ux430ux434ux430ux447ux438-ux43cux430ux448ux438ux43dux43dux43eux433ux43e-ux43eux431ux443ux447ux435ux43dux438ux44f}}

\begin{enumerate}
\def\labelenumi{\arabic{enumi}.}
\tightlist
\item
  \textbf{Обучение с учителем:}

  \begin{itemize}
  \tightlist
  \item
    \textbf{Регрессия} --- прогнозирование численных данных (например,
    прибыль).
  \item
    \textbf{Классификация} --- предсказание категорий (кликнет ли
    пользователь на рекламу).
  \end{itemize}
\item
  \textbf{Обучение без учителя:}

  \begin{itemize}
  \tightlist
  \item
    \textbf{Кластеризация} --- разделение объектов на группы.
  \item
    \textbf{Оценка плотности} --- поиск аномалий.
  \item
    \textbf{Понижение размерности} --- упрощение данных.
  \end{itemize}
\item
  \textbf{Обучение с подкреплением} --- оптимизация действий для
  максимизации награды (например, управление беспилотным автомобилем).
\end{enumerate}

\hypertarget{ux43aux43bux44eux447ux435ux432ux44bux435-ux44dux442ux430ux43fux44b-ux440ux430ux431ux43eux442ux44b}{%
\subsection*{Ключевые этапы
работы}\label{ux43aux43bux44eux447ux435ux432ux44bux435-ux44dux442ux430ux43fux44b-ux440ux430ux431ux43eux442ux44b}}

\begin{enumerate}
\def\labelenumi{\arabic{enumi}.}
\tightlist
\item
  Постановка задачи.
\item
  Выделение признаков.
\item
  Формирование выборки.
\item
  Выбор функционала ошибки.
\item
  Предобработка данных.
\item
  Обучение модели.
\item
  Оценка качества.
\end{enumerate}

\hypertarget{ux43fux440ux438ux43cux435ux440ux44b-ux43fux440ux438ux43cux435ux43dux435ux43dux438ux44f}{%
\subsection*{Примеры
применения}\label{ux43fux440ux438ux43cux435ux440ux44b-ux43fux440ux438ux43cux435ux43dux435ux43dux438ux44f}}

\hypertarget{ux43fux440ux435ux438ux43cux443ux449ux435ux441ux442ux432ux430}{%
\subsubsection*{Преимущества:}\label{ux43fux440ux435ux438ux43cux443ux449ux435ux441ux442ux432ux430}}

\begin{itemize}
\tightlist
\item
  Выявление скрытых закономерностей.
\item
  Прогнозирование сложных зависимостей.
\item
  Автоматизация трудоёмких задач (например, классификация овощей по
  качеству).
\end{itemize}

\hypertarget{ux443ux441ux43bux43eux432ux438ux44f-ux443ux441ux43fux435ux448ux43dux43eux433ux43e-ux438ux441ux43fux43eux43bux44cux437ux43eux432ux430ux43dux438ux44f}{%
\subsubsection*{Условия успешного
использования:}\label{ux443ux441ux43bux43eux432ux438ux44f-ux443ux441ux43fux435ux448ux43dux43eux433ux43e-ux438ux441ux43fux43eux43bux44cux437ux43eux432ux430ux43dux438ux44f}}

\begin{itemize}
\tightlist
\item
  Достаточный объём данных.
\item
  Сложность зависимости между признаками и целевой переменной.
\item
  Невозможность построения точных формул.
\end{itemize}

\hypertarget{ux43bux438ux43dux435ux439ux43dux44bux435-ux43cux43eux434ux435ux43bux438}{%
\subsection*{Линейные
модели}\label{ux43bux438ux43dux435ux439ux43dux44bux435-ux43cux43eux434ux435ux43bux438}}

Линейные модели представляют собой один из базовых инструментов
машинного обучения. Они описываются формулой:

\begin{equation}
a(x) = w_0 + \sum_{j=1}^{d} w_j x_j,
\end{equation}


где \(w_j\)--- веса модели, \(w_0\)--- свободный коэффициент,
\(x_j\)--- признаки объекта.

Такая форма позволяет выразить прогноз как скалярное произведение
вектора признаков и весов.

\hypertarget{ux43fux440ux435ux438ux43cux443ux449ux435ux441ux442ux432ux430-ux43bux438ux43dux435ux439ux43dux44bux445-ux43cux43eux434ux435ux43bux435ux439}{%
\subsubsection*{Преимущества линейных
моделей:}\label{ux43fux440ux435ux438ux43cux443ux449ux435ux441ux442ux432ux430-ux43bux438ux43dux435ux439ux43dux44bux445-ux43cux43eux434ux435ux43bux435ux439}}

\begin{enumerate}
\def\labelenumi{\arabic{enumi}.}
\tightlist
\item
  \textbf{Простота обучения} --- вычисления происходят быстро.
\item
  \textbf{Устойчивость} --- ограниченное количество параметров снижает
  риск переобучения.
\item
  \textbf{Интерпретируемость} --- вклад каждого признака легко
  анализировать.
\end{enumerate}

Однако линейные модели предполагают линейную зависимость между
признаками и целевой переменной, что ограничивает их применимость для
сложных зависимостей.

\hypertarget{ux43eux431ux43bux430ux441ux442ux438-ux43fux440ux438ux43cux435ux43dux435ux43dux438ux44f-ux438-ux43eux431ux440ux430ux431ux43eux442ux43aux430-ux434ux430ux43dux43dux44bux445}{%
\subsection*{Области применения и обработка
данных}\label{ux43eux431ux43bux430ux441ux442ux438-ux43fux440ux438ux43cux435ux43dux435ux43dux438ux44f-ux438-ux43eux431ux440ux430ux431ux43eux442ux43aux430-ux434ux430ux43dux43dux44bux445}}

\hypertarget{ux43aux430ux442ux435ux433ux43eux440ux438ux430ux43bux44cux43dux44bux435-ux43fux440ux438ux437ux43dux430ux43aux438}{%
\subsubsection*{1. Категориальные
признаки:}\label{ux43aux430ux442ux435ux433ux43eux440ux438ux430ux43bux44cux43dux44bux435-ux43fux440ux438ux437ux43dux430ux43aux438}}

\begin{itemize}
\tightlist
\item
  Кодируются через \textbf{one-hot представление} --- преобразование
  признака в несколько бинарных переменных.\\
  Пример: если признак принимает значения \{A, B, C\}, он преобразуется
  в три бинарных признака (один из которых равен единице, остальные ---
  нулю).
\end{itemize}

\hypertarget{ux440ux430ux431ux43eux442ux430-ux441-ux442ux435ux43aux441ux442ux430ux43cux438}{%
\subsubsection*{2. Работа с
текстами:}\label{ux440ux430ux431ux43eux442ux430-ux441-ux442ux435ux43aux441ux442ux430ux43cux438}}

\begin{itemize}
\tightlist
\item
  Метод ``мешок слов'' (bag of words): текст представляется как набор
  уникальных слов, где каждый признак отражает частоту появления
  конкретного слова в тексте.
\end{itemize}

\hypertarget{ux431ux438ux43dux430ux440ux438ux437ux430ux446ux438ux44f-ux447ux438ux441ux43bux43eux432ux44bux445-ux43fux440ux438ux437ux43dux430ux43aux43eux432}{%
\subsubsection*{3. Бинаризация числовых
признаков:}\label{ux431ux438ux43dux430ux440ux438ux437ux430ux446ux438ux44f-ux447ux438ux441ux43bux43eux432ux44bux445-ux43fux440ux438ux437ux43dux430ux43aux43eux432}}

\begin{itemize}
\tightlist
\item
  Разбиение данных на интервалы. Каждый интервал представляется бинарным
  признаком. Это позволяет учитывать нелинейные зависимости.
\end{itemize}

\hypertarget{ux43cux435ux442ux440ux438ux43aux438-ux43eux446ux435ux43dux43aux438-ux43aux430ux447ux435ux441ux442ux432ux430-ux43cux43eux434ux435ux43bux435ux439}{%
\subsection*{Метрики оценки качества
моделей}\label{ux43cux435ux442ux440ux438ux43aux438-ux43eux446ux435ux43dux43aux438-ux43aux430ux447ux435ux441ux442ux432ux430-ux43cux43eux434ux435ux43bux435ux439}}

\hypertarget{ux441ux440ux435ux434ux43dux435ux43aux432ux430ux434ux440ux430ux442ux438ux447ux43dux430ux44f-ux43eux448ux438ux431ux43aux430-mse}{%
\subsubsection*{1. Среднеквадратичная ошибка
(MSE):}\label{ux441ux440ux435ux434ux43dux435ux43aux432ux430ux434ux440ux430ux442ux438ux447ux43dux430ux44f-ux43eux448ux438ux431ux43aux430-mse}}

\[
MSE = \frac{1}{\ell} \sum_{i=1}^{\ell} (a(x_i) - y_i)^2.
\]

\begin{itemize}
\tightlist
\item
  Подходит для оптимизации благодаря дифференцируемости.
\item
  Чувствительна к выбросам.
\item
  Часто используется корень MSE (RMSE) для сохранения единиц измерения.
\end{itemize}

\hypertarget{ux441ux440ux435ux434ux43dux44fux44f-ux430ux431ux441ux43eux43bux44eux442ux43dux430ux44f-ux43eux448ux438ux431ux43aux430-mae}{%
\subsubsection*{2. Средняя абсолютная ошибка
(MAE):}\label{ux441ux440ux435ux434ux43dux44fux44f-ux430ux431ux441ux43eux43bux44eux442ux43dux430ux44f-ux43eux448ux438ux431ux43aux430-mae}}

\[
MAE = \frac{1}{\ell} \sum_{i=1}^{\ell} |a(x_i) - y_i|.
\]

\begin{itemize}
\tightlist
\item
  Устойчивее к выбросам, чем MSE.
\item
  Труднее оптимизируется.
\end{itemize}

\hypertarget{ux43aux43eux44dux444ux444ux438ux446ux438ux435ux43dux442-ux434ux435ux442ux435ux440ux43cux438ux43dux430ux446ux438ux438-r2}{%
\subsubsection*{\texorpdfstring{3. Коэффициент детерминации
(\(R^2\)):}{3. Коэффициент детерминации (R\^{}2):}}\label{ux43aux43eux44dux444ux444ux438ux446ux438ux435ux43dux442-ux434ux435ux442ux435ux440ux43cux438ux43dux430ux446ux438ux438-r2}}

\[
R^2 = 1 - \frac{\sum_{i=1}^{\ell} (a(x_i) - y_i)^2}{\sum_{i=1}^{\ell} (y_i - \bar{y})^2}.
\]

\begin{itemize}
\tightlist
\item
  \(R^2\), близкий к 1, указывает на хорошую модель, а близкий к 0 ---
  на слабую.
\end{itemize}

\hypertarget{ux433ux438ux431ux440ux438ux434ux43dux44bux435-ux444ux443ux43dux43aux446ux438ux438-ux43fux43eux442ux435ux440ux44c}{%
\subsubsection*{4. Гибридные функции
потерь:}\label{ux433ux438ux431ux440ux438ux434ux43dux44bux435-ux444ux443ux43dux43aux446ux438ux438-ux43fux43eux442ux435ux440ux44c}}

\begin{itemize}
\tightlist
\item
  \textbf{Huber Loss} --- комбинирует свойства MSE и MAE:
\end{itemize}

\[
L_\delta(y, a) = 
\begin{cases} 
\frac{1}{2}(y - a)^2, & \text{если } |y - a| < \delta, \\
\delta (|y - a| - \frac{1}{2}\delta), & \text{иначе}.
\end{cases}
\]

\begin{itemize}
\tightlist
\item
  \textbf{Log-Cosh Loss} --- мягче штрафует большие ошибки:
\end{itemize}

\[
L(y, a) = \log(\cosh(y - a)).
\]

\hypertarget{ux441ux440ux435ux434ux43dux435ux43aux432ux430ux434ux440ux430ux442ux438ux447ux43dux430ux44f-ux43bux43eux433ux430ux440ux438ux444ux43cux438ux447ux435ux441ux43aux430ux44f-ux43eux448ux438ux431ux43aux430-msle}{%
\subsubsection*{5. Среднеквадратичная логарифмическая ошибка
(MSLE):}\label{ux441ux440ux435ux434ux43dux435ux43aux432ux430ux434ux440ux430ux442ux438ux447ux43dux430ux44f-ux43bux43eux433ux430ux440ux438ux444ux43cux438ux447ux435ux441ux43aux430ux44f-ux43eux448ux438ux431ux43aux430-msle}}

\[
MSLE = \frac{1}{\ell} \sum_{i=1}^{\ell} (\log(a(x_i) + 1) - \log(y_i + 1))^2.
\]

\begin{itemize}
\tightlist
\item
  Применяется для неотрицательных целевых переменных.
\end{itemize}

\hypertarget{ux43fux440ux43eux446ux435ux43dux442ux43dux44bux435-ux43eux448ux438ux431ux43aux438-mape-ux438-smape}{%
\subsubsection*{6. Процентные ошибки (MAPE и
SMAPE):}\label{ux43fux440ux43eux446ux435ux43dux442ux43dux44bux435-ux43eux448ux438ux431ux43aux438-mape-ux438-smape}}

\begin{itemize}
\tightlist
\item
  \textbf{MAPE} измеряет среднюю абсолютную процентную ошибку:
\end{itemize}

\[
MAPE = \frac{1}{\ell} \sum_{i=1}^{\ell} \left|\frac{y_i - a(x_i)}{y_i}\right|.
\]

\begin{itemize}
\tightlist
\item
  \textbf{SMAPE} нормализует ошибку на среднее значение прогнозов и
  ответов:
\end{itemize}

\[
SMAPE = \frac{|y - a|}{(|y| + |a|)/2}.
\]

\hypertarget{ux43aux432ux430ux43dux442ux438ux43bux44cux43dux430ux44f-ux444ux443ux43dux43aux446ux438ux44f-ux43fux43eux442ux435ux440ux44c}{%
\subsubsection*{7. Квантильная функция
потерь:}\label{ux43aux432ux430ux43dux442ux438ux43bux44cux43dux430ux44f-ux444ux443ux43dux43aux446ux438ux44f-ux43fux43eux442ux435ux440ux44c}}

\begin{itemize}
\tightlist
\item
  Используется для задач с разной ценностью ошибок:
\end{itemize}

\[
Q(a, X) = \sum_{i=1}^{\ell} \rho_\tau(y_i - a(x_i)),
\]

где \(\rho_\tau(z)\) акцентирует внимание на занижениях или завышениях.

\hypertarget{ux43fux435ux440ux435ux43eux431ux443ux447ux435ux43dux438ux435}{%
\section{Переобучение}\label{ux43fux435ux440ux435ux43eux431ux443ux447ux435ux43dux438ux435}}

Переобучение возникает, когда модель чрезмерно подстраивается под
обучающую выборку, теряя способность обобщать закономерности на новые
данные. Это часто проявляется при использовании слишком сложных моделей,
например, с большим количеством признаков или их высокими степенями.

\textbf{Пример:} модель с признаками \(\{x, x^2, \dots, x^{15}\}\)может
точно описать обучающую выборку, но плохо справляется с предсказанием
новых данных, так как учитывает случайный шум.

\hypertarget{ux43cux435ux442ux43eux434ux44b-ux431ux43eux440ux44cux431ux44b-ux441-ux43fux435ux440ux435ux43eux431ux443ux447ux435ux43dux438ux435ux43c}{%
\subsubsection*{Методы борьбы с
переобучением:}\label{ux43cux435ux442ux43eux434ux44b-ux431ux43eux440ux44cux431ux44b-ux441-ux43fux435ux440ux435ux43eux431ux443ux447ux435ux43dux438ux435ux43c}}

\begin{itemize}
\tightlist
\item
  \textbf{Сужение класса моделей:} выбираются модели меньшей сложности.
\item
  \textbf{Регуляризация:} вводится штраф за большие значения
  коэффициентов модели, чтобы уменьшить влияние случайного шума.
\end{itemize}

\hypertarget{ux43eux446ux435ux43dux438ux432ux430ux43dux438ux435-ux43aux430ux447ux435ux441ux442ux432ux430-ux43cux43eux434ux435ux43bux435ux439}{%
\section{Оценивание качества
моделей}\label{ux43eux446ux435ux43dux438ux432ux430ux43dux438ux435-ux43aux430ux447ux435ux441ux442ux432ux430-ux43cux43eux434ux435ux43bux435ux439}}

Качество моделей оценивается с помощью отложенной выборки.

\hypertarget{ux434ux430ux43dux43dux44bux435-ux434ux435ux43bux44fux442ux441ux44f-ux43dux430}{%
\subsubsection*{Данные делятся
на:}\label{ux434ux430ux43dux43dux44bux435-ux434ux435ux43bux44fux442ux441ux44f-ux43dux430}}

\begin{enumerate}
\def\labelenumi{\arabic{enumi}.}
\tightlist
\item
  \textbf{Обучающую выборку} --- используется для настройки модели.
\item
  \textbf{Контрольную выборку} --- предназначена для проверки качества
  модели.
\end{enumerate}

\hypertarget{ux43aux440ux43eux441ux441-ux432ux430ux43bux438ux434ux430ux446ux438ux44f}{%
\subsubsection*{Кросс-валидация:}\label{ux43aux440ux43eux441ux441-ux432ux430ux43bux438ux434ux430ux446ux438ux44f}}

Для повышения надёжности используется метод кросс-валидации: 
\begin{enumerate}
\def\labelenumi{\arabic{enumi}.}
\tightlist
\item 
Данные делятся на \(k\) частей (блоков). 
\item 
Модель обучается на \(k-1\) блоках и
тестируется на оставшемся. 
\item 
Процесс повторяется \(k\) раз, каждый раз
меняя тестовый блок. 
\item 
Итоговое качество модели вычисляется как среднее
по всем итерациям.
\end{enumerate}

\hypertarget{ux43fux43eux441ux43bux435-ux43aux440ux43eux441ux441-ux432ux430ux43bux438ux434ux430ux446ux438ux438-ux432ux43eux437ux43cux43eux436ux43dux44b-ux434ux432ux430-ux43fux43eux434ux445ux43eux434ux430}{%
\subsubsection*{После кросс-валидации возможны два
подхода:}\label{ux43fux43eux441ux43bux435-ux43aux440ux43eux441ux441-ux432ux430ux43bux438ux434ux430ux446ux438ux438-ux432ux43eux437ux43cux43eux436ux43dux44b-ux434ux432ux430-ux43fux43eux434ux445ux43eux434ux430}}

\begin{enumerate}
\def\labelenumi{\arabic{enumi}.}
\tightlist
\item
  \textbf{Обучение на всей выборке:} модель обучается на всех доступных
  данных, чтобы извлечь максимум информации.
\item
  \textbf{Композиция моделей:} объединение моделей из каждого шага
  кросс-валидации (например, усреднение прогнозов в регрессии).
\end{enumerate}

\hypertarget{ux43eux431ux443ux447ux435ux43dux438ux435-ux43bux438ux43dux435ux439ux43dux43eux439-ux440ux435ux433ux440ux435ux441ux441ux438ux438}{%
\section{Обучение линейной
регрессии}\label{ux43eux431ux443ux447ux435ux43dux438ux435-ux43bux438ux43dux435ux439ux43dux43eux439-ux440ux435ux433ux440ux435ux441ux441ux438ux438}}

Линейная регрессия обучается путём минимизации среднеквадратичной ошибки
(MSE):

\[
\frac{1}{\ell} \sum_{i=1}^\ell (\langle w, x_i \rangle - y_i)^2 \to \min_w,
\]

где \(X\) --- матрица признаков, \(y\) --- вектор целевых переменных,
\(w\) --- вектор параметров модели.

\hypertarget{ux440ux435ux448ux435ux43dux438ux435-ux437ux430ux434ux430ux447ux438}{%
\subsubsection*{Решение
задачи:}\label{ux440ux435ux448ux435ux43dux438ux435-ux437ux430ux434ux430ux447ux438}}

Представлено в матричном виде:

\[
w = (X^TX)^{-1}X^Ty.
\]

\hypertarget{ux43fux440ux43eux431ux43bux435ux43cux44b-ux430ux43dux430ux43bux438ux442ux438ux447ux435ux441ux43aux43eux433ux43e-ux440ux435ux448ux435ux43dux438ux44f}{%
\subsubsection*{Проблемы аналитического
решения:}\label{ux43fux440ux43eux431ux43bux435ux43cux44b-ux430ux43dux430ux43bux438ux442ux438ux447ux435ux441ux43aux43eux433ux43e-ux440ux435ux448ux435ux43dux438ux44f}}

\begin{enumerate}
\def\labelenumi{\arabic{enumi}.}
\tightlist
\item
  \textbf{Сложность вычислений:} обращение матрицы требует
  \(O(d^3)\) операций, что непрактично для больших \(d\).
\item
  \textbf{Плохо обусловленная матрица:} если \(X^TX\) близка к
  вырожденной, результат становится нестабильным.\\
  \textbf{Решение:} регуляризация или численные методы.
\end{enumerate}

\hypertarget{ux433ux440ux430ux434ux438ux435ux43dux442ux43dux44bux439-ux441ux43fux443ux441ux43a}{%
\section{Градиентный
спуск}\label{ux433ux440ux430ux434ux438ux435ux43dux442ux43dux44bux439-ux441ux43fux443ux441ux43a}}

Градиентный спуск --- это итеративный метод оптимизации, подходящий для
задач, где аналитическое решение недоступно.

\hypertarget{ux448ux430ux433ux438-ux430ux43bux433ux43eux440ux438ux442ux43cux430}{%
\subsubsection*{Шаги
алгоритма:}\label{ux448ux430ux433ux438-ux430ux43bux433ux43eux440ux438ux442ux43cux430}}

\[
w^{(k)} = w^{(k-1)} - \eta_k \nabla Q(w^{(k-1)}),
\]

где  \(Q(w)\) --- функция ошибки, \(\eta_k\) --- длина шага, 
\(\nabla Q\) --- градиент функции ошибки.

\hypertarget{ux43eux441ux43dux43eux432ux43dux44bux435-ux441ux432ux43eux439ux441ux442ux432ux430-ux433ux440ux430ux434ux438ux435ux43dux442ux430}{%
\subsubsection*{Основные свойства
градиента:}\label{ux43eux441ux43dux43eux432ux43dux44bux435-ux441ux432ux43eux439ux441ux442ux432ux430-ux433ux440ux430ux434ux438ux435ux43dux442ux430}}

\begin{enumerate}
\def\labelenumi{\arabic{enumi}.}
\tightlist
\item
  Градиент показывает направление наибольшего роста функции.
\item
  Антиградиент \(-\nabla Q\) указывает направление наискорейшего
  уменьшения.
\item
  Градиент ортогонален линиям уровня функции.
\end{enumerate}

\hypertarget{ux432ux44bux431ux43eux440-ux434ux43bux438ux43dux44b-ux448ux430ux433ux430-eta_k}{%
\subsubsection*{\texorpdfstring{Выбор длины шага
\(\eta_k\):}{Выбор длины шага (\textbackslash eta\_k):}}\label{ux432ux44bux431ux43eux440-ux434ux43bux438ux43dux44b-ux448ux430ux433ux430-eta_k}}

\begin{itemize}
\tightlist
\item
  \textbf{Слишком большой шаг} может привести к нестабильности
  (перепрыгиванию минимума).
\item
  \textbf{Слишком маленький шаг} замедляет сходимость.
\item
  На практике используются:

  \begin{itemize}
  \tightlist
  \item
    Адаптивные шаги,
  \item
    Убывающий шаг, например \(\eta_k = \frac{1}{k}\).
  \end{itemize}
\end{itemize}

\hypertarget{ux43cux435ux442ux43eux434ux44b-ux43eux446ux435ux43dux438ux432ux430ux43dux438ux44f-ux433ux440ux430ux434ux438ux435ux43dux442ux430}{%
\section{Методы оценивания
градиента}\label{ux43cux435ux442ux43eux434ux44b-ux43eux446ux435ux43dux438ux432ux430ux43dux438ux44f-ux433ux440ux430ux434ux438ux435ux43dux442ux430}}

Для ускорения работы вместо полного градиента для всей выборки:

\[
\nabla Q(w) = \frac{1}{\ell} \sum_{i=1}^\ell \nabla q_i(w),
\]

применяются приближённые методы:

\hypertarget{ux441ux442ux43eux445ux430ux441ux442ux438ux447ux435ux441ux43aux438ux439-ux433ux440ux430ux434ux438ux435ux43dux442ux43dux44bux439-ux441ux43fux443ux441ux43a-sgd}{%
\subsubsection*{1. Стохастический градиентный спуск
(SGD):}\label{ux441ux442ux43eux445ux430ux441ux442ux438ux447ux435ux441ux43aux438ux439-ux433ux440ux430ux434ux438ux435ux43dux442ux43dux44bux439-ux441ux43fux443ux441ux43a-sgd}}

\begin{itemize}
\tightlist
\item
  Градиент вычисляется на одном случайно выбранном объекте:
\end{itemize}

\[
\nabla Q(w) \approx \nabla q_i(w).
\]

\begin{itemize}
\tightlist
\item
  \textbf{Преимущества:} высокая скорость итерации.
\item
  \textbf{Недостатки:} менее точен, требует убывающего шага для
  сходимости.
\end{itemize}

\hypertarget{mini-batch-ux433ux440ux430ux434ux438ux435ux43dux442ux43dux44bux439-ux441ux43fux443ux441ux43a}{%
\subsubsection*{2. Mini-batch градиентный
спуск:}\label{mini-batch-ux433ux440ux430ux434ux438ux435ux43dux442ux43dux44bux439-ux441ux43fux443ux441ux43a}}

\begin{itemize}
\tightlist
\item
  Использует подвыборки данных (пакеты) для оценки градиента:
\end{itemize}

\[
\nabla Q(w) \approx \frac{1}{n} \sum_{j=1}^n \nabla q_{i_j}(w),
\]

где \(n\) --- размер пакета.

\hypertarget{stochastic-average-gradient-sag}{%
\subsubsection*{3. Stochastic Average Gradient
(SAG):}\label{stochastic-average-gradient-sag}}

\begin{itemize}
\tightlist
\item
  Хранит все вычисленные ранее градиенты.
\item
  Обновляет градиент только для одного объекта на каждом шаге.
\item
  \textbf{Преимущество:} ускоряет сходимость.
\item
  \textbf{Недостаток:} требует памяти для хранения всех градиентов.
\end{itemize}

\hypertarget{ux43cux43eux434ux438ux444ux438ux43aux430ux446ux438ux438-ux433ux440ux430ux434ux438ux435ux43dux442ux43dux43eux433ux43e-ux441ux43fux443ux441ux43aux430}{%
\paragraph{Модификации градиентного
спуска}\label{ux43cux43eux434ux438ux444ux438ux43aux430ux446ux438ux438-ux433ux440ux430ux434ux438ux435ux43dux442ux43dux43eux433ux43e-ux441ux43fux443ux441ux43aux430}}

\begin{enumerate}
\def\labelenumi{\arabic{enumi}.}
\tightlist
\item
  \textbf{Метод инерции (Momentum):}

  \begin{itemize}
  \tightlist
  \item
    Уменьшает осцилляции, усредняя градиенты нескольких итераций:
  \end{itemize}


\[
h_k = \alpha h_{k-1} + \eta_k \nabla Q(w^{(k-1)}),
\]

где \(\alpha\) --- коэффициент инерции. 

\begin{itemize}
  \tightlist
  \item
Позволяет быстрее находить минимум.
\end{itemize}
\end{enumerate}

\begin{enumerate}
\def\labelenumi{\arabic{enumi}.}
\setcounter{enumi}{1}
\tightlist
\item
  \textbf{AdaGrad:}

  \begin{itemize}
  \tightlist
  \item
    Адаптирует длину шага для каждого параметра:
  \end{itemize}
\end{enumerate}

\[
w_j^{(k)} = w_j^{(k-1)} - \frac{\eta}{\sqrt{G_{k,j} + \epsilon}} \nabla_j Q(w^{(k-1)})
\]

где \(G_{k,j}\) --- накопленные квадраты градиентов, а
\(\epsilon\)предотвращает деление на ноль.

\begin{enumerate}
\def\labelenumi{\arabic{enumi}.}
\setcounter{enumi}{2}
\tightlist
\item
  \textbf{RMSprop:}

  \begin{itemize}
  \tightlist
  \item
    Вводит экспоненциальное затухание градиентов, чтобы шаги не
    замедлялись слишком сильно:
  \end{itemize}
\end{enumerate}

\[
G_{k,j} = \alpha G_{k-1,j} + (1 - \alpha)(\nabla_j Q(w^{(k-1)}))^2.
\]

\begin{enumerate}
\def\labelenumi{\arabic{enumi}.}
\setcounter{enumi}{3}
\tightlist
\item
  \textbf{Adam:}

  \begin{itemize}
  \tightlist
  \item
    Комбинирует идеи Momentum и RMSprop, обеспечивая быструю и
    стабильную сходимость.
  \end{itemize}
\end{enumerate}

\hypertarget{ux43bux438ux43dux435ux439ux43dux44bux435-ux43cux43eux434ux435ux43bux438-ux43aux43bux430ux441ux441ux438ux444ux438ux43aux430ux446ux438ux438}{%
\subsection*{Линейные модели
классификации}\label{ux43bux438ux43dux435ux439ux43dux44bux435-ux43cux43eux434ux435ux43bux438-ux43aux43bux430ux441ux441ux438ux444ux438ux43aux430ux446ux438ux438}}

\hypertarget{ux444ux43eux440ux43cux430ux43bux438ux437ux430ux446ux438ux44f}{%
\subsubsection*{Формализация}\label{ux444ux43eux440ux43cux430ux43bux438ux437ux430ux446ux438ux44f}}

Линейная модель классификации для задачи бинарной классификации
определяется следующим образом:

\[
a(x) = \text{sign} \left( \langle w, x \rangle + w_0 \right) = \text{sign} \left( \sum_{j=1}^d w_j x_j + w_0 \right),
\]

где \(w \) --- вектор весов, \(w_0 \) --- сдвиг (bias). Модель разделяет
пространство объектов гиперплоскостью, относит полупространства к
положительному или отрицательному классу.

Если в признаках добавляется константа \(x_{d+1} = 1 \), сдвиг \(w_0\) становится ненужным, и модель принимает вид:

\[
a(x) = \text{sign} \langle w, x \rangle.
\]

\hypertarget{ux43eux431ux443ux447ux435ux43dux438ux435-ux43bux438ux43dux435ux439ux43dux44bux445-ux43aux43bux430ux441ux441ux438ux444ux438ux43aux430ux442ux43eux440ux43eux432}{%
\subsection*{Обучение линейных
классификаторов}\label{ux43eux431ux443ux447ux435ux43dux438ux435-ux43bux438ux43dux435ux439ux43dux44bux445-ux43aux43bux430ux441ux441ux438ux444ux438ux43aux430ux442ux43eux440ux43eux432}}

\hypertarget{ux43eux441ux43dux43eux432ux43dux43eux439-ux444ux443ux43dux43aux446ux438ux43eux43dux430ux43b}{%
\subsubsection*{Основной
функционал}\label{ux43eux441ux43dux43eux432ux43dux43eux439-ux444ux443ux43dux43aux446ux438ux43eux43dux430ux43b}}

Для обучения используется доля неправильных ответов:

\[
Q(a, X) = \frac{1}{\ell} \sum_{i=1}^\ell \left[ \text{sign} \langle w, x_i \rangle \neq y_i \right] \to \min_w.
\]

Этот функционал является дискретным и сложен для минимизации. Для
решения задачи вводят \textbf{отступы}:

\[
M_i = y_i \langle w, x_i \rangle.
\]

Функционал можно переписать через отступы:

\[
Q(a, X) = \frac{1}{\ell} \sum_{i=1}^\ell [M_i < 0].
\]

Для минимизации используют \textbf{гладкие верхние оценки} на функцию
потерь \(L(M) = {[}M \textless{} 0{]} \), например: 
\begin{itemize}
\item
Логистическая:
\(L(M) = \log(1 + e^{-M})\),
\item
Метод опорных векторов:
\(L(M) = \max(0, 1 - M)\), 
\item
Персептрон: \(L(M) = \max(0, -M)\).
\end{itemize}
\hypertarget{ux43cux435ux442ux440ux438ux43aux438-ux43aux430ux447ux435ux441ux442ux432ux430-ux43aux43bux430ux441ux441ux438ux444ux438ux43aux430ux446ux438ux438}{%
\subsection*{Метрики качества
классификации}\label{ux43cux435ux442ux440ux438ux43aux438-ux43aux430ux447ux435ux441ux442ux432ux430-ux43aux43bux430ux441ux441ux438ux444ux438ux43aux430ux446ux438ux438}}

\hypertarget{ux434ux43eux43bux44f-ux43fux440ux430ux432ux438ux43bux44cux43dux44bux445-ux43eux442ux432ux435ux442ux43eux432}{%
\subsubsection*{1. Доля правильных
ответов}\label{ux434ux43eux43bux44f-ux43fux440ux430ux432ux438ux43bux44cux43dux44bux445-ux43eux442ux432ux435ux442ux43eux432}}

Определяется как:

\[
\text{accuracy} = \frac{1}{\ell} \sum_{i=1}^\ell [a(x_i) = y_i].
\]

Недостаток: игнорирует несбалансированность классов. Для анализа вводят
\textbf{относительное уменьшение ошибки}.

\hypertarget{ux43cux430ux442ux440ux438ux446ux430-ux43eux448ux438ux431ux43eux43a}{%
\subsubsection*{2. Матрица
ошибок}\label{ux43cux430ux442ux440ux438ux446ux430-ux43eux448ux438ux431ux43eux43a}}

Объекты классификации разделяют на категории: \(\text{TP}\)(True
Positive), \(\text{FP}\)(False Positive), \(\text{FN}\)(False Negative),
\(\text{TN}\)(True Negative).

Отсюда вводят: 
\begin{itemize}
\tightlist
\item
Точность (precision):
\end{itemize}

\[
\text{precision} = \frac{\text{TP}}{\text{TP} + \text{FP}},
\]

\begin{itemize}
\tightlist
\item
  Полноту (recall):
\end{itemize}

\[
\text{recall} = \frac{\text{TP}}{\text{TP} + \text{FN}}.
\]

Также используется \textbf{F-мера} --- гармоническое среднее точности и
полноты:

\[
F = \frac{2 \cdot \text{precision} \cdot \text{recall}}{\text{precision} + \text{recall}}.
\]

\hypertarget{roc-ux43aux440ux438ux432ux430ux44f-ux438-auc-roc}{%
\subsubsection*{3. ROC-кривая и
AUC-ROC}\label{roc-ux43aux440ux438ux432ux430ux44f-ux438-auc-roc}}

ROC-кривая строится в координатах: - False Positive Rate (FPR):

\[
\text{FPR} = \frac{\text{FP}}{\text{FP} + \text{TN}},
\]

\begin{itemize}
\tightlist
\item
  True Positive Rate (TPR):
\end{itemize}

\[
\text{TPR} = \frac{\text{TP}}{\text{TP} + \text{FN}}.
\]

Площадь под ROC-кривой \( \text{AUC-ROC} \) измеряет обобщённое
качество алгоритма. В задачах кредитного скоринга применяется
\textbf{индекс Джини}:

\[
\text{Gini} = 2 \cdot \text{AUC} - 1.
\]

\hypertarget{precision-recall-ux43aux440ux438ux432ux430ux44f-ux438-auc-pr}{%
\subsubsection*{4. Precision-Recall кривая и
AUC-PR}\label{precision-recall-ux43aux440ux438ux432ux430ux44f-ux438-auc-pr}}

Используется для несбалансированных классов. Precision-Recall кривая
оценивает точность и полноту, строя их зависимость от порога
классификации. Площадь под PR-кривой
\(( \text{AUC-PR} )\) аппроксимируется через среднюю точность:

\[
\text{AUC-PR} = \frac{1}{\ell_+} \sum_{k=1}^\ell [y_k = 1] \cdot \text{precision@k}.
\]

\hypertarget{ux43bux43eux433ux438ux441ux442ux438ux447ux435ux441ux43aux430ux44f-ux440ux435ux433ux440ux435ux441ux441ux438ux44f}{%
\subsection*{Логистическая
регрессия}\label{ux43bux43eux433ux438ux441ux442ux438ux447ux435ux441ux43aux430ux44f-ux440ux435ux433ux440ux435ux441ux441ux438ux44f}}

\hypertarget{ux43eux446ux435ux43dux438ux432ux430ux43dux438ux435-ux432ux435ux440ux43eux44fux442ux43dux43eux441ux442ux435ux439}{%
\subsubsection*{Оценивание
вероятностей}\label{ux43eux446ux435ux43dux438ux432ux430ux43dux438ux435-ux432ux435ux440ux43eux44fux442ux43dux43eux441ux442ux435ux439}}

Логистическая регрессия строится на основе логистической функции потерь.
Основная задача --- корректно оценивать вероятность принадлежности
объекта к классу \(+1 \).

Пусть вероятность принадлежности объекта \(x \) к классу \(+1
\) обозначена как \(p(y = +1 \textbar{} x) \). Для алгоритма \(b(x) \),
возвращающего вероятности, минимизация функционала должна приводить к:

\[
\arg\min_{b \in \mathbb{R}} \frac{1}{n} \sum_{i=1}^n L(y_i, b) \approx p(y = +1 | x).
\]

Пример: логарифмическая функция потерь позволяет корректно оценивать
вероятности. Если объект \$x \$встречается \$ n \$ раз, матожидание
ошибки имеет вид:

\[
\arg\min_{b \in \mathbb{R}} \left( \frac{k}{n} L(1, b) + \frac{n-k}{n} L(-1, b) \right) = \frac{k}{n},
\]

где \( k \) --- количество объектов класса \( +1 \).

\hypertarget{ux43fux440ux430ux432ux434ux43eux43fux43eux434ux43eux431ux438ux435-ux438-ux43bux43eux433ux438ux441ux442ux438ux447ux435ux441ux43aux438ux435-ux43fux43eux442ux435ux440ux438}{%
\subsubsection*{Правдоподобие и логистические
потери}\label{ux43fux440ux430ux432ux434ux43eux43fux43eux434ux43eux431ux438ux435-ux438-ux43bux43eux433ux438ux441ux442ux438ux447ux435ux441ux43aux438ux435-ux43fux43eux442ux435ux440ux438}}

Логистическая функция потерь \( \text{log-loss} \) основана на
максимизации правдоподобия выборки:

\[
Q(a, X) = \prod_{i=1}^\ell b(x_i)^{[y_i=+1]} (1 - b(x_i))^{[y_i=-1]}.
\]

Логарифм правдоподобия удобен для оптимизации:

\[
-\sum_{i=1}^\ell \left( [y_i = +1] \log b(x_i) + [y_i = -1] \log(1 - b(x_i)) \right) \to \min.
\]

Эта функция потерь оптимизирует правдоподобие выборки и обеспечивает
корректные оценки вероятности.

\hypertarget{ux43bux43eux433ux438ux441ux442ux438ux447ux435ux441ux43aux430ux44f-ux440ux435ux433ux440ux435ux441ux441ux438ux44f-1}{%
\subsubsection*{Логистическая
регрессия}\label{ux43bux43eux433ux438ux441ux442ux438ux447ux435ux441ux43aux430ux44f-ux440ux435ux433ux440ux435ux441ux441ux438ux44f-1}}

Логистическая регрессия преобразует линейную модель через сигмоидную
функцию:

\[
b(x) = \sigma(\langle w, x \rangle), \quad \sigma(z) = \frac{1}{1 + \exp(-z)}.
\]

Скалярное произведение \( \langle w, x \rangle \) интерпретируется как
логарифм отношения вероятностей классов \( \text{log-odds} \):

\[
\langle w, x \rangle = \log \frac{p(y = +1 | x)}{p(y = -1 | x)}.
\]

Логистические потери можно выразить следующим образом:

\[
-\sum_{i=1}^\ell \log(1 + \exp(-y_i \langle w, x_i \rangle)).
\]

Оптимизация этого функционала определяет модель логистической регрессии.

\hypertarget{ux43cux435ux442ux43eux434-ux43eux43fux43eux440ux43dux44bux445-ux432ux435ux43aux442ux43eux440ux43eux432}{%
\subsection*{Метод опорных
векторов}\label{ux43cux435ux442ux43eux434-ux43eux43fux43eux440ux43dux44bux445-ux432ux435ux43aux442ux43eux440ux43eux432}}

Метод опорных векторов (SVM) основан на максимизации разделяющего зазора
между классами.

\hypertarget{ux440ux430ux437ux434ux435ux43bux438ux43cux44bux439-ux441ux43bux443ux447ux430ux439}{%
\subsubsection*{Разделимый
случай}\label{ux440ux430ux437ux434ux435ux43bux438ux43cux44bux439-ux441ux43bux443ux447ux430ux439}}

Если выборка линейно разделима, задача SVM записывается как:

\[
\begin{aligned}
&\frac{1}{2} \|w\|^2 \to \min_{w, b}, \\
&y_i (\langle w, x_i \rangle + b) \geq 1, \quad i = 1, \ldots, \ell.
\end{aligned}
\]

Максимизация ширины разделяющей полосы улучшает обобщающую способность
классификатора.

\hypertarget{ux43dux435ux440ux430ux437ux434ux435ux43bux438ux43cux44bux439-ux441ux43bux443ux447ux430ux439}{%
\subsubsection*{Неразделимый
случай}\label{ux43dux435ux440ux430ux437ux434ux435ux43bux438ux43cux44bux439-ux441ux43bux443ux447ux430ux439}}

Если выборка не разделима, вводятся штрафы за нарушения \( \xi_i
\textgreater{} 0 \):

\[
y_i (\langle w, x_i \rangle + b) \geq 1 - \xi_i, \quad \xi_i \geq 0.
\]

Оптимизационная задача принимает вид:

\[
\frac{1}{2} \|w\|^2 + C \sum_{i=1}^\ell \xi_i \to \min_{w, b, \xi}.
\]

Параметр \(C\)регулирует баланс между максимизацией отступа и
минимизацией штрафов.

\hypertarget{ux441ux432ux435ux434ux435ux43dux438ux435-ux43a-ux431ux435ux437ux443ux441ux43bux43eux432ux43dux43eux439-ux437ux430ux434ux430ux447ux435}{%
\subsubsection*{Сведение к безусловной
задаче}\label{ux441ux432ux435ux434ux435ux43dux438ux435-ux43a-ux431ux435ux437ux443ux441ux43bux43eux432ux43dux43eux439-ux437ux430ux434ux430ux447ux435}}

Штрафы \(\xi_i\) можно выразить через кусочно-линейную функцию потерь \text{hinge loss}:

\[
\xi_i = \max(0, 1 - y_i (\langle w, x_i \rangle + b)).
\]

Тогда функционал становится:

\[
\frac{1}{2} \|w\|^2 + C \sum_{i=1}^\ell \max(0, 1 - y_i (\langle w, x_i \rangle + b)) \to \min_{w, b}.
\]

Этот функционал использует верхнюю оценку на долю ошибок с добавлением
регуляризации.

\hypertarget{ux43cux43dux43eux433ux43eux43aux43bux430ux441ux441ux43eux432ux430ux44f-ux43aux43bux430ux441ux441ux438ux444ux438ux43aux430ux446ux438ux44f}{%
\subsection*{Многоклассовая
классификация}\label{ux43cux43dux43eux433ux43eux43aux43bux430ux441ux441ux43eux432ux430ux44f-ux43aux43bux430ux441ux441ux438ux444ux438ux43aux430ux446ux438ux44f}}

\hypertarget{ux441ux432ux435ux434ux435ux43dux438ux435-ux43a-ux441ux435ux440ux438ux438-ux431ux438ux43dux430ux440ux43dux44bux445-ux437ux430ux434ux430ux447}{%
\subsubsection*{Сведение к серии бинарных
задач}\label{ux441ux432ux435ux434ux435ux43dux438ux435-ux43a-ux441ux435ux440ux438ux438-ux431ux438ux43dux430ux440ux43dux44bux445-ux437ux430ux434ux430ux447}}

\begin{enumerate}
\def\labelenumi{\arabic{enumi}.}
\tightlist
\item
  \textbf{Один против всех (One-vs-All):}

  \begin{itemize}
  \tightlist
  \item
    Обучается \(K\) классификаторов \(b_1(x), \ldots, b_K(x)\), каждый
    из которых отличает один класс от остальных.
  \item
    Итоговый классификатор выбирает класс с наибольшей уверенностью:
  \end{itemize}
\end{enumerate}

\[
a(x) = \arg\max_{k \in \{1, \ldots, K\}} b_k(x).
\]

\begin{enumerate}
\def\labelenumi{\arabic{enumi}.}
\setcounter{enumi}{1}
\tightlist
\item
  \textbf{Все против всех (All-vs-All):}

  \begin{itemize}
  \tightlist
  \item
    Обучается \(\binom{K}{2}\) классификаторов \(a_{ij}(x)\), каждый из
    которых различает пару классов \(i\) и \(j\).
  \item
    Класс определяется большинством голосов.
  \end{itemize}
\end{enumerate}

\hypertarget{ux43cux43dux43eux433ux43eux43aux43bux430ux441ux441ux43eux432ux430ux44f-ux43bux43eux433ux438ux441ux442ux438ux447ux435ux441ux43aux430ux44f-ux440ux435ux433ux440ux435ux441ux441ux438ux44f}{%
\subsubsection*{Многоклассовая логистическая
регрессия}\label{ux43cux43dux43eux433ux43eux43aux43bux430ux441ux441ux43eux432ux430ux44f-ux43bux43eux433ux438ux441ux442ux438ux447ux435ux441ux43aux430ux44f-ux440ux435ux433ux440ux435ux441ux441ux438ux44f}}

\begin{itemize}
\tightlist
\item
  Логистическая регрессия обобщается с использованием оператора
  \textbf{SoftMax}:
\end{itemize}

\[
\text{SoftMax}(z_1, \ldots, z_K) = \left( \frac{\exp(z_1)}{\sum_{k=1}^K \exp(z_k)}, \ldots, \frac{\exp(z_K)}{\sum_{k=1}^K \exp(z_k)} \right).
\]

\begin{itemize}
\tightlist
\item
  Вероятность принадлежности класса \(k\):
\end{itemize}

\[
P(y = k | x, w) = \frac{\exp(\langle w_k, x \rangle + w_{0k})}{\sum_{j=1}^K \exp(\langle w_j, x \rangle + w_{0j})}.
\]

\begin{itemize}
\tightlist
\item
  Обучение происходит методом максимального правдоподобия.
\end{itemize}

\hypertarget{ux43cux43dux43eux433ux43eux43aux43bux430ux441ux441ux43eux432ux44bux439-ux43cux435ux442ux43eux434-ux43eux43fux43eux440ux43dux44bux445-ux432ux435ux43aux442ux43eux440ux43eux432-svm}{%
\subsubsection*{Многоклассовый метод опорных векторов
(SVM)}\label{ux43cux43dux43eux433ux43eux43aux43bux430ux441ux441ux43eux432ux44bux439-ux43cux435ux442ux43eux434-ux43eux43fux43eux440ux43dux44bux445-ux432ux435ux43aux442ux43eux440ux43eux432-svm}}

\begin{enumerate}
\def\labelenumi{\arabic{enumi}.}
\tightlist
\item
  \textbf{Функция потерь:}
\end{enumerate}

\[
\max_k \{ \langle w_k, x \rangle + 1 - [k = y(x)] \} - \langle w_{y(x)}, x \rangle.
\]

Потери равны нулю, если разрыв между правильным и другими классами
превышает 1.

\begin{enumerate}
\def\labelenumi{\arabic{enumi}.}
\setcounter{enumi}{1}
\tightlist
\item
  \textbf{Линейно разделимая выборка:}

  \begin{itemize}
  \tightlist
  \item
    Оптимизация максимального отступа (аналог бинарного SVM):
  \end{itemize}
\end{enumerate}

\[
\frac{1}{2} \|W\|^2 \to \min_W, \quad \langle w_{y_i}, x_i \rangle + [y_i = k] - \langle w_k, x_i \rangle > 1.
\]

\begin{enumerate}
\def\labelenumi{\arabic{enumi}.}
\setcounter{enumi}{2}
\tightlist
\item
  \textbf{Неразделимая выборка:}

  \begin{itemize}
  \tightlist
  \item
    Добавляются штрафы \(\xi_i > 0\):
  \end{itemize}
\end{enumerate}

\[
\frac{1}{2} \|W\|^2 + C \sum_{i=1}^\ell \xi_i \to \min_{W, \xi}.
\]

\hypertarget{ux43cux435ux442ux440ux438ux43aux438-ux43aux430ux447ux435ux441ux442ux432ux430-ux43cux43dux43eux433ux43eux43aux43bux430ux441ux441ux43eux432ux43eux439-ux43aux43bux430ux441ux441ux438ux444ux438ux43aux430ux446ux438ux438}{%
\subsubsection*{Метрики качества многоклассовой
классификации}\label{ux43cux435ux442ux440ux438ux43aux438-ux43aux430ux447ux435ux441ux442ux432ux430-ux43cux43dux43eux433ux43eux43aux43bux430ux441ux441ux43eux432ux43eux439-ux43aux43bux430ux441ux441ux438ux444ux438ux43aux430ux446ux438ux438}}

Используются микро- и макро-усреднения: 

- \textbf{Микро-усреднение:}
усреднение характеристик (TP, FP и т.д.) по всем классам. 

- \textbf{Макро-усреднение:} усреднение метрик для каждого класса.

Пример: точность в микро- и макро-усреднении: 
\begin{itemize}
\item
  Микро:
\end{itemize}
\[
\text{precision} = \frac{\text{TP}}{\text{TP} + \text{FP}}.
\]

\begin{itemize}
\tightlist
\item
  Макро:
\end{itemize}

\[
\text{precision} = \frac{1}{K} \sum_{k=1}^K \frac{\text{TP}_k}{\text{TP}_k + \text{FP}_k}.
\]

\hypertarget{ux43aux43bux430ux441ux441ux438ux444ux438ux43aux430ux446ux438ux44f-ux441-ux43fux435ux440ux435ux441ux435ux43aux430ux44eux449ux438ux43cux438ux441ux44f-ux43aux43bux430ux441ux441ux430ux43cux438-multi-label}{%
\subsection*{Классификация с пересекающимися классами
(Multi-label)}\label{ux43aux43bux430ux441ux441ux438ux444ux438ux43aux430ux446ux438ux44f-ux441-ux43fux435ux440ux435ux441ux435ux43aux430ux44eux449ux438ux43cux438ux441ux44f-ux43aux43bux430ux441ux441ux430ux43cux438-multi-label}}

\hypertarget{ux43fux43eux434ux445ux43eux434ux44b-ux43a-ux440ux435ux448ux435ux43dux438ux44e}{%
\subsubsection*{Подходы к
решению}\label{ux43fux43eux434ux445ux43eux434ux44b-ux43a-ux440ux435ux448ux435ux43dux438ux44e}}

\begin{enumerate}
\def\labelenumi{\arabic{enumi}.}
\tightlist
\item
  \textbf{Независимая классификация (Binary Relevance):}

  \begin{itemize}
  \tightlist
  \item
    Обучается \(K\) независимых классификаторов
    \(a_1(x), \ldots, a_K(x)\).
  \end{itemize}
\item
  \textbf{Стекинг классификаторов:}

  \begin{itemize}
  \tightlist
  \item
    Прогнозы базовых моделей \(b_k(x)\) используются в качестве
    признаков для обучения новых моделей \(a_k(x')\).
  \end{itemize}
\item
  \textbf{Трансформация пространства ответов:}

  \begin{itemize}
  \tightlist
  \item
    Используется сингулярное разложение матрицы ответов \(Y\):
  \end{itemize}
\end{enumerate}

\[
Y = U \Sigma V^\top.
\]

\begin{itemize}
\tightlist
\item
  Сжимаются ответы до \(M\)-мерного пространства.
\end{itemize}

\hypertarget{ux43cux435ux442ux440ux438ux43aux438-ux43aux430ux447ux435ux441ux442ux432ux430}{%
\subsubsection*{Метрики
качества}\label{ux43cux435ux442ux440ux438ux43aux438-ux43aux430ux447ux435ux441ux442ux432ux430}}

\begin{enumerate}
\def\labelenumi{\arabic{enumi}.}
\tightlist
\item
  \textbf{Хэмминговая метрика:}
\end{enumerate}

\[
\text{hamming}(a, X) = \frac{1}{\ell} \sum_{i=1}^\ell \frac{|Y_i \setminus Z_i| + |Z_i \setminus Y_i|}{K}.
\]

\begin{enumerate}
\def\labelenumi{\arabic{enumi}.}
\setcounter{enumi}{1}
\tightlist
\item
  Обобщённые метрики:

  \begin{itemize}
  \tightlist
  \item
    Точность:
  \end{itemize}


\[
\text{accuracy}(a, X) = \frac{1}{\ell} \sum_{i=1}^\ell \frac{|Y_i \cap Z_i|}{|Y_i \cup Z_i|}.
\]

\begin{itemize}
\tightlist
\item
  Полнота:
\end{itemize}

\[
\text{recall}(a, X) = \frac{1}{\ell} \sum_{i=1}^\ell \frac{|Y_i \cap Z_i|}{|Y_i|}.
\]

\hypertarget{ux43aux430ux442ux435ux433ux43eux440ux438ux430ux43bux44cux43dux44bux435-ux43fux440ux438ux437ux43dux430ux43aux438-1}{%
\subsection*{Категориальные
признаки}\label{ux43aux430ux442ux435ux433ux43eux440ux438ux430ux43bux44cux43dux44bux435-ux43fux440ux438ux437ux43dux430ux43aux438-1}}

\hypertarget{ux43fux43eux434ux445ux43eux434ux44b-ux43a-ux43eux431ux440ux430ux431ux43eux442ux43aux435}{%
\subsubsection*{Подходы к
обработке}\label{ux43fux43eux434ux445ux43eux434ux44b-ux43a-ux43eux431ux440ux430ux431ux43eux442ux43aux435}}

\begin{enumerate}
\def\labelenumi{\arabic{enumi}.}
\tightlist
\item
  \textbf{Бинарное кодирование (One-Hot Encoding):}

    Создаются \(n\) индикаторов \(g_j(x) = [f(x) = u_j]\).
    
\item
  \textbf{Бинарное кодирование с хэшированием:}

    Значения признаков переводятся в \(B\)-размерное пространство с
    помощью хэш-функции.
    
\item
  \textbf{Счётчики:}

  \begin{itemize}
  \tightlist
  \item
    Считаются статистики по выборке:
  \end{itemize}
\end{enumerate}

\[
\text{counts}(u, X) = \sum_{(x, y) \in X} [f(x) = u], \quad \text{successes}_k(u, X) = \sum_{(x, y) \in X} [f(x) = u][y = k].
\]

    Признаки оценивают вероятность \(p(y = k | f(x))\).



\hypertarget{ux440ux435ux448ux430ux44eux449ux438ux435-ux434ux435ux440ux435ux432ux44cux44f-ux438-ux440ux430ux43dux434ux43eux43cux43dux44bux439-ux43bux435ux441}{%
\subsection*{Решающие деревья и Рандомный
лес}\label{ux440ux435ux448ux430ux44eux449ux438ux435-ux434ux435ux440ux435ux432ux44cux44f-ux438-ux440ux430ux43dux434ux43eux43cux43dux44bux439-ux43bux435ux441}}

\hypertarget{ux447ux442ux43e-ux442ux430ux43aux43eux435-ux440ux435ux448ux430ux44eux449ux438ux435-ux434ux435ux440ux435ux432ux44cux44f}{%
\subsubsection*{Что такое решающие
деревья?}\label{ux447ux442ux43e-ux442ux430ux43aux43eux435-ux440ux435ux448ux430ux44eux449ux438ux435-ux434ux435ux440ux435ux432ux44cux44f}}

Решающее дерево --- это алгоритм машинного обучения, который
представляет собой иерархическую структуру, использующую
последовательные условия для разделения данных на группы. Каждый узел
дерева отвечает за проверку условия на одном из признаков, а листовые
узлы дают прогноз или решение.

Решающие деревья популярны благодаря их интерпретируемости и способности
работать как с непрерывными, так и с категориальными данными. Они
используются для задач классификации и регрессии.

\hypertarget{ux43fux440ux438ux43dux446ux438ux43f-ux440ux430ux431ux43eux442ux44b}{%
\subsubsection*{Принцип
работы}\label{ux43fux440ux438ux43dux446ux438ux43f-ux440ux430ux431ux43eux442ux44b}}

\begin{enumerate}
\def\labelenumi{\arabic{enumi}.}
\tightlist
\item
  \textbf{Разделение данных:}

  \begin{itemize}
  \tightlist
  \item
    Дерево строится путем рекурсивного разделения пространства признаков
    на основе условий. Каждое разделение основывается на максимизации
    ``чистоты'' полученных подмножеств данных.
  \end{itemize}
\item
  \textbf{Оценка качества разделения:}

  \begin{itemize}
  \tightlist
  \item
    Для классификации используется \textbf{критерий чистоты}, например:

    \begin{itemize}
    \tightlist
    \item
      \textbf{Индекс Джини:}
    \end{itemize}
  \end{itemize}
\end{enumerate}

\[
G = 1 - \sum_{k=1}^K p_k^2,
\]

где \(p_k\) --- доля объектов класса \(k\). 

\begin{itemize}
    \tightlist
    \item
\textbf{Энтропия:}
\end{itemize}
\end{enumerate}

\[
H = - \sum_{k=1}^K p_k \log(p_k).
\]

\begin{itemize}
\tightlist
\item
  Для регрессии используется \textbf{дисперсия} или
  \textbf{среднеквадратичная ошибка}:
\end{itemize}

\[ 
MSE = \frac{1}{n} \sum_{i=1}^n (y_i - \hat{y})^2.
\]

\begin{enumerate}
\def\labelenumi{\arabic{enumi}.}
\setcounter{enumi}{3}
\tightlist
\item
  \textbf{Остановка роста дерева:}

  \begin{itemize}
  \tightlist
  \item
    Процесс деления останавливается, когда:

    \begin{itemize}
    \tightlist
    \item
      Достигнута минимальная глубина дерева.
    \item
      В узле осталось меньше определенного числа объектов.
    \item
      Новое разделение не улучшает качество модели.
    \end{itemize}
  \end{itemize}
\end{enumerate}

\hypertarget{ux43fux440ux435ux438ux43cux443ux449ux435ux441ux442ux432ux430-ux440ux435ux448ux430ux44eux449ux438ux445-ux434ux435ux440ux435ux432ux44cux435ux432}{%
\subsubsection*{Преимущества решающих
деревьев}\label{ux43fux440ux435ux438ux43cux443ux449ux435ux441ux442ux432ux430-ux440ux435ux448ux430ux44eux449ux438ux445-ux434ux435ux440ux435ux432ux44cux435ux432}}

\begin{enumerate}
\def\labelenumi{\arabic{enumi}.}
\tightlist
\item
  Интерпретируемость: легко понять логику модели.
\item
  Поддержка смешанных данных (числовых и категориальных).
\item
  Нет необходимости в масштабировании признаков.
\item
  Хорошо работает с пропусками данных.
\end{enumerate}

\hypertarget{ux43dux435ux434ux43eux441ux442ux430ux442ux43aux438-ux440ux435ux448ux430ux44eux449ux438ux445-ux434ux435ux440ux435ux432ux44cux435ux432}{%
\subsubsection*{Недостатки решающих
деревьев}\label{ux43dux435ux434ux43eux441ux442ux430ux442ux43aux438-ux440ux435ux448ux430ux44eux449ux438ux445-ux434ux435ux440ux435ux432ux44cux435ux432}}

\begin{enumerate}
\def\labelenumi{\arabic{enumi}.}
\tightlist
\item
  Склонность к переобучению (если дерево слишком глубокое).
\item
  Чувствительность к небольшим изменениям данных (нестабильность).
\item
  Ограниченная способность к обобщению для сложных задач.
\end{enumerate}

\hypertarget{ux43eux431ux440ux435ux437ux43aux430-pruning}{%
\subsubsection*{Обрезка
(Pruning)}\label{ux43eux431ux440ux435ux437ux43aux430-pruning}}

\begin{itemize}
\tightlist
\item
  \textbf{Предварительная обрезка (Pre-pruning):}

  \begin{itemize}
  \tightlist
  \item
    Ограничение роста дерева (например, максимальная глубина или
    минимальное количество объектов в узле).
  \end{itemize}
\item
  \textbf{Пост-обрезка (Post-pruning):}

  \begin{itemize}
  \tightlist
  \item
    Удаление лишних ветвей после полного построения дерева, если они не
    улучшают качество.
  \end{itemize}
\end{itemize}

\hypertarget{ux440ux435ux433ux443ux43bux44fux440ux438ux437ux430ux446ux438ux44f}{%
\subsubsection*{Регуляризация}\label{ux440ux435ux433ux443ux43bux44fux440ux438ux437ux430ux446ux438ux44f}}

\begin{itemize}
\tightlist
\item
  Включение штрафа за сложность модели:

  \begin{itemize}
  \tightlist
  \item
    Ограничение глубины.
  \item
    Установка минимального числа объектов в узле.
  \item
    Ограничение минимального прироста качества для разделения.
  \end{itemize}
\end{itemize}

\hypertarget{ux447ux442ux43e-ux442ux430ux43aux43eux435-ux440ux430ux43dux434ux43eux43cux43dux44bux439-ux43bux435ux441}{%
\subsubsection*{Что такое рандомный
лес?}\label{ux447ux442ux43e-ux442ux430ux43aux43eux435-ux440ux430ux43dux434ux43eux43cux43dux44bux439-ux43bux435ux441}}

Рандомный лес --- это ансамблевый метод, объединяющий множество решающих
деревьев. Итоговый прогноз строится на основе голосования деревьев (для
классификации) или усреднения их прогнозов (для регрессии).

\hypertarget{ux43eux441ux43dux43eux432ux43dux44bux435-ux438ux434ux435ux438-ux440ux430ux43dux434ux43eux43cux43dux43eux433ux43e-ux43bux435ux441ux430}{%
\subsubsection*{Основные идеи рандомного
леса}\label{ux43eux441ux43dux43eux432ux43dux44bux435-ux438ux434ux435ux438-ux440ux430ux43dux434ux43eux43cux43dux43eux433ux43e-ux43bux435ux441ux430}}

\begin{enumerate}
\def\labelenumi{\arabic{enumi}.}
\tightlist
\item
  \textbf{Бутстреппинг (Bagging):}

  \begin{itemize}
  \tightlist
  \item
    Каждое дерево обучается на случайной подвыборке исходных данных.
  \item
    Это снижает корреляцию между деревьями и увеличивает стабильность
    модели.
  \end{itemize}
\item
  \textbf{Случайный выбор признаков:}

  \begin{itemize}
  \tightlist
  \item
    Для каждого узла дерева случайным образом выбирается подмножество
    признаков, на которых строится разделение.
  \item
    Это предотвращает доминирование сильных признаков.
  \end{itemize}
\item
  \textbf{Ансамблевое обучение:}

  \begin{itemize}
  \tightlist
  \item
    Итоговый прогноз получается путем объединения прогнозов всех
    деревьев:

    \begin{itemize}
    \tightlist
    \item
      Для классификации: \textbf{большинство голосов}.
    \item
      Для регрессии: \textbf{среднее значение}.
    \end{itemize}
  \end{itemize}
\end{enumerate}

\hypertarget{ux43fux440ux435ux438ux43cux443ux449ux435ux441ux442ux432ux430-ux440ux430ux43dux434ux43eux43cux43dux43eux433ux43e-ux43bux435ux441ux430}{%
\subsubsection*{Преимущества рандомного
леса}\label{ux43fux440ux435ux438ux43cux443ux449ux435ux441ux442ux432ux430-ux440ux430ux43dux434ux43eux43cux43dux43eux433ux43e-ux43bux435ux441ux430}}

\begin{enumerate}
\def\labelenumi{\arabic{enumi}.}
\tightlist
\item
  Высокая точность на многих задачах.
\item
  Стабильность и устойчивость к шуму.
\item
  Устойчивость к переобучению благодаря ансамблевому подходу.
\item
  Поддержка работы с большим числом признаков.
\end{enumerate}

\hypertarget{ux43dux435ux434ux43eux441ux442ux430ux442ux43aux438-ux440ux430ux43dux434ux43eux43cux43dux43eux433ux43e-ux43bux435ux441ux430}{%
\subsubsection*{Недостатки рандомного
леса}\label{ux43dux435ux434ux43eux441ux442ux430ux442ux43aux438-ux440ux430ux43dux434ux43eux43cux43dux43eux433ux43e-ux43bux435ux441ux430}}

\begin{enumerate}
\def\labelenumi{\arabic{enumi}.}
\tightlist
\item
  Потеря интерпретируемости (модель становится «черным ящиком»).
\item
  Высокая вычислительная сложность для больших наборов данных.
\item
  Модель может быть менее эффективной, если признаки несут слабую
  информацию.
\end{enumerate}

\hypertarget{ux432ux430ux436ux43dux43eux441ux442ux44c-ux43fux440ux438ux437ux43dux430ux43aux43eux432}{%
\subsubsection*{Важность
признаков}\label{ux432ux430ux436ux43dux43eux441ux442ux44c-ux43fux440ux438ux437ux43dux430ux43aux43eux432}}

\begin{itemize}
\tightlist
\item
  В рандомном лесе можно оценить значимость признаков:

  \begin{itemize}
  \tightlist
  \item
    \textbf{Decrease in Impurity:} оценка по уменьшению критерия
    (например, Джини) при разделении.
  \item
    \textbf{Permutation Importance:} оценка по изменению качества модели
    при случайной перестановке значений признака.
  \end{itemize}
\end{itemize}

\hypertarget{ux43dux430ux441ux442ux440ux43eux439ux43aux430-ux433ux438ux43fux435ux440ux43fux430ux440ux430ux43cux435ux442ux440ux43eux432}{%
\subsubsection*{Настройка
гиперпараметров}\label{ux43dux430ux441ux442ux440ux43eux439ux43aux430-ux433ux438ux43fux435ux440ux43fux430ux440ux430ux43cux435ux442ux440ux43eux432}}

\hypertarget{ux434ux43bux44f-ux440ux435ux448ux430ux44eux449ux438ux445-ux434ux435ux440ux435ux432ux44cux435ux432}{%
\paragraph{Для решающих
деревьев:}\label{ux434ux43bux44f-ux440ux435ux448ux430ux44eux449ux438ux445-ux434ux435ux440ux435ux432ux44cux435ux432}}

\begin{itemize}
\tightlist
\item
  \texttt{max\_depth}: максимальная глубина дерева.
\item
  \texttt{min\_samples\_split}: минимальное количество объектов для
  разделения узла.
\item
  \texttt{min\_samples\_leaf}: минимальное количество объектов в листе.
\item
  \texttt{criterion}: функция оценки качества (например, \texttt{gini}
  или \texttt{entropy}).
\end{itemize}

\hypertarget{ux434ux43bux44f-ux440ux430ux43dux434ux43eux43cux43dux43eux433ux43e-ux43bux435ux441ux430}{%
\paragraph{Для рандомного
леса:}\label{ux434ux43bux44f-ux440ux430ux43dux434ux43eux43cux43dux43eux433ux43e-ux43bux435ux441ux430}}

\begin{itemize}
\tightlist
\item
  \texttt{n\_estimators}: количество деревьев в ансамбле.
\item
  \texttt{max\_features}: число признаков для выбора при разделении.
\item
  \texttt{bootstrap}: использование бутстреппинга (да/нет).
\item
  \texttt{max\_depth}, \texttt{min\_samples\_split},
  \texttt{min\_samples\_leaf}: параметры деревьев внутри леса.
\end{itemize}

\hypertarget{ux43fux440ux438ux43cux435ux440ux44b-ux43fux440ux438ux43cux435ux43dux435ux43dux438ux44f-1}{%
\subsubsection*{Примеры
применения}\label{ux43fux440ux438ux43cux435ux440ux44b-ux43fux440ux438ux43cux435ux43dux435ux43dux438ux44f-1}}

\begin{enumerate}
\def\labelenumi{\arabic{enumi}.}
\tightlist
\item
  \textbf{Классификация:}

  \begin{itemize}
  \tightlist
  \item
    Определение заболеваний на основе медицинских данных.
  \item
    Классификация изображений и текстов.
  \item
    Предсказание оттока клиентов.
  \end{itemize}
\item
  \textbf{Регрессия:}

  \begin{itemize}
  \tightlist
  \item
    Прогнозирование цен на недвижимость.
  \item
    Оценка риска в кредитовании.
  \item
    Прогнозирование спроса на товары.
  \end{itemize}
\item
  \textbf{Выявление важности признаков:}

  \begin{itemize}
  \tightlist
  \item
    Анализ ключевых факторов, влияющих на целевую переменную.
  \item
    Обнаружение нерелевантных или избыточных признаков.
  \end{itemize}
\end{enumerate}

\hypertarget{ux43fux440ux430ux43aux442ux438ux447ux435ux441ux43aux438ux435-ux440ux435ux43aux43eux43cux435ux43dux434ux430ux446ux438ux438}{%
\subsubsection*{Практические
рекомендации}\label{ux43fux440ux430ux43aux442ux438ux447ux435ux441ux43aux438ux435-ux440ux435ux43aux43eux43cux435ux43dux434ux430ux446ux438ux438}}

\begin{enumerate}
\def\labelenumi{\arabic{enumi}.}
\tightlist
\item
  \textbf{Для решающих деревьев:}

  \begin{itemize}
  \tightlist
  \item
    Начинайте с небольших глубин и увеличивайте их при необходимости.
  \item
    Используйте кросс-валидацию для выбора оптимальных гиперпараметров.
  \end{itemize}
\item
  \textbf{Для рандомного леса:}

  \begin{itemize}
  \tightlist
  \item
    Увеличивайте количество деревьев \(n\_estimators\), пока качество не
    стабилизируется.
  \item
    Используйте бутстреппинг и ограничение числа признаков
    \(max\_features\) для повышения качества.
  \end{itemize}
\item
  \textbf{Предобработка данных:}

  \begin{itemize}
  \tightlist
  \item
    Заполните пропуски или используйте подходы, устойчивые к пропускам.
  \item
    Преобразуйте категориальные данные (например, с помощью one-hot
    encoding).
  \end{itemize}
\end{enumerate}

\hypertarget{ux431ux443ux441ux442ux438ux43dux433}{%
\subsubsection*{Бустинг}\label{ux431ux443ux441ux442ux438ux43dux433}}

Бустинг --- это ансамблевый метод машинного обучения, который строит модель
путем последовательного добавления слабых моделей (обычно решающих
деревьев) с целью улучшения качества предсказаний.

\hypertarget{ux432ux432ux435ux434ux435ux43dux438ux435-ux432-ux431ux443ux441ux442ux438ux43dux433}{%
\subsubsection*{Введение в
бустинг}\label{ux432ux432ux435ux434ux435ux43dux438ux435-ux432-ux431ux443ux441ux442ux438ux43dux433}}

\hypertarget{ux43eux441ux43dux43eux432ux43dux430ux44f-ux438ux434ux435ux44f-ux431ux443ux441ux442ux438ux43dux433ux430}{%
\subsubsection*{Основная идея
бустинга}\label{ux43eux441ux43dux43eux432ux43dux430ux44f-ux438ux434ux435ux44f-ux431ux443ux441ux442ux438ux43dux433ux430}}

\begin{enumerate}
\def\labelenumi{\arabic{enumi}.}
\tightlist
\item
  \textbf{Слабые модели:} В качестве базовых используются слабые
  алгоритмы, которые имеют качество чуть лучше случайного угадывания
  (например, мелкие решающие деревья).
\item
  \textbf{Итеративное обучение:} Новая модель настраивается так, чтобы
  исправлять ошибки предыдущих моделей.
\item
  \textbf{Ансамбль:} Итоговый прогноз получается путем взвешенного
  суммирования всех слабых моделей.
\end{enumerate}

\hypertarget{ux442ux438ux43fux44b-ux431ux443ux441ux442ux438ux43dux433ux430}{%
\subsubsection*{Типы
бустинга}\label{ux442ux438ux43fux44b-ux431ux443ux441ux442ux438ux43dux433ux430}}

\hypertarget{ux430ux434ux430ux43fux442ux438ux432ux43dux44bux439-ux431ux443ux441ux442ux438ux43dux433-adaboost}{%
\subsubsection*{Адаптивный бустинг
(AdaBoost)}\label{ux430ux434ux430ux43fux442ux438ux432ux43dux44bux439-ux431ux443ux441ux442ux438ux43dux433-adaboost}}

\begin{enumerate}
\def\labelenumi{\arabic{enumi}.}
\tightlist
\item
  \textbf{Алгоритм:}

  \begin{itemize}
  \tightlist
  \item
    На каждой итерации \(t\) строится слабый классификатор \(h_t(x)\).
  \item
    Ошибки предыдущих моделей усиливаются за счет изменения весов
    объектов.
  \end{itemize}
\item
  \textbf{Формулы:}

  \begin{itemize}
  \tightlist
  \item
    Начальные веса объектов:
  \end{itemize}
\end{enumerate}

\[ 
     w_i^{(1)} = \frac{1}{N}, 
\]

где \(N\) --- количество объектов. - Ошибка текущей модели:

\[ 
     \varepsilon_t = \frac{\sum_{i=1}^N w_i^{(t)} \cdot [y_i \neq h_t(x_i)]}{\sum_{i=1}^N w_i^{(t)}}.
\]

\begin{itemize}
\tightlist
\item
  Коэффициент модели:
\end{itemize}

\[ 
     \alpha_t = \frac{1}{2} \ln\left(\frac{1 - \varepsilon_t}{\varepsilon_t}\right).
\]

\begin{itemize}
\tightlist
\item
  Обновление весов:
\end{itemize}

\[
     w_i^{(t+1)} = w_i^{(t)} \cdot \exp\left(-\alpha_t \cdot y_i \cdot h_t(x_i)\right).
\]

\begin{enumerate}
\def\labelenumi{\arabic{enumi}.}
\setcounter{enumi}{2}
\tightlist
\item
  Итоговый ансамбль:
\end{enumerate}

\[
   H(x) = \text{sign} \left( \sum_{t=1}^T \alpha_t \cdot h_t(x) \right).
\]

\hypertarget{ux433ux440ux430ux434ux438ux435ux43dux442ux43dux44bux439-ux431ux443ux441ux442ux438ux43dux433-gradient-boosting}{%
\subsubsection*{Градиентный бустинг (Gradient
Boosting)}\label{ux433ux440ux430ux434ux438ux435ux43dux442ux43dux44bux439-ux431ux443ux441ux442ux438ux43dux433-gradient-boosting}}

\begin{enumerate}
\def\labelenumi{\arabic{enumi}.}
\tightlist
\item
  \textbf{Идея:}

  \begin{itemize}
  \tightlist
  \item
    Каждая новая модель минимизирует остатки (ошибки) предыдущих
    моделей.
  \item
    Для регрессии минимизируется разница между истинными значениями и
    предсказаниями.
  \item
    Для классификации минимизируется логарифмическая функция потерь.
  \end{itemize}
\item
  \textbf{Формулы:}

  \begin{itemize}
  \tightlist
  \item
    Остатки на \(t\)-й итерации:
  \end{itemize}
\end{enumerate}

\[
     r_i^{(t)} = -\frac{\partial L(y_i, F(x_i))}{\partial F(x_i)},
\]


 где \(L(y, F(x))\) — функция потерь.


\begin{itemize}
\tightlist
\item
  Новый прогноз: \[
  F_{t+1}(x) = F_t(x) + \eta \cdot h_t(x),
  \]
\end{itemize}

где \(\eta\) --- шаг обучения, \(h_t(x)\) --- новое дерево.

\begin{enumerate}
\def\labelenumi{\arabic{enumi}.}
\setcounter{enumi}{2}
\tightlist
\item
  Итоговый ансамбль:
\end{enumerate}

\[
   F_T(x) = \sum_{t=1}^T \eta \cdot h_t(x).
\]

\hypertarget{ux441ux442ux43eux445ux430ux441ux442ux438ux447ux435ux441ux43aux438ux439-ux433ux440ux430ux434ux438ux435ux43dux442ux43dux44bux439-ux431ux443ux441ux442ux438ux43dux433-stochastic-gradient-boosting}{%
\subsubsection*{Стохастический градиентный бустинг (Stochastic Gradient
Boosting)}\label{ux441ux442ux43eux445ux430ux441ux442ux438ux447ux435ux441ux43aux438ux439-ux433ux440ux430ux434ux438ux435ux43dux442ux43dux44bux439-ux431ux443ux441ux442ux438ux43dux433-stochastic-gradient-boosting}}

\begin{itemize}
\tightlist
\item
  Отличие от классического градиентного бустинга:

  \begin{itemize}
  \tightlist
  \item
    На каждой итерации используется случайная подвыборка объектов
    (например, 50\%-70\%).
  \item
    Это снижает переобучение и увеличивает обобщающую способность.
  \end{itemize}
\end{itemize}

\hypertarget{ux43aux430ux442ux435ux433ux43eux440ux438ux430ux43bux44cux43dux44bux439-ux431ux443ux441ux442ux438ux43dux433-catboost}{%
\subsubsection*{Категориальный бустинг
(CatBoost)}\label{ux43aux430ux442ux435ux433ux43eux440ux438ux430ux43bux44cux43dux44bux439-ux431ux443ux441ux442ux438ux43dux433-catboost}}

\begin{enumerate}
\def\labelenumi{\arabic{enumi}.}
\tightlist
\item
  \textbf{Особенности:}

  \begin{itemize}
  \tightlist
  \item
    Адаптирован для категориальных признаков.
  \item
    Использует кодирование категориальных признаков на основе счётчиков
    с учётом информации о целевой переменной.
  \end{itemize}
\item
  \textbf{Обучение:}

  \begin{itemize}
  \tightlist
  \item
    Упрощает обработку категориальных признаков.
  \item
    Улучшает качество предсказания за счёт оптимизированных алгоритмов.
  \end{itemize}
\end{enumerate}

\hypertarget{ux43aux43eux43cux43fux43eux43dux435ux43dux442ux44b-ux431ux443ux441ux442ux438ux43dux433ux430}{%
\subsubsection*{Компоненты
бустинга}\label{ux43aux43eux43cux43fux43eux43dux435ux43dux442ux44b-ux431ux443ux441ux442ux438ux43dux433ux430}}

\hypertarget{ux431ux430ux437ux43eux432ux44bux435-ux43cux43eux434ux435ux43bux438-ux441ux43bux430ux431ux44bux435-ux43cux43eux434ux435ux43bux438}{%
\subsubsection*{Базовые модели (слабые
модели)}\label{ux431ux430ux437ux43eux432ux44bux435-ux43cux43eux434ux435ux43bux438-ux441ux43bux430ux431ux44bux435-ux43cux43eux434ux435ux43bux438}}

\begin{enumerate}
\def\labelenumi{\arabic{enumi}.}
\tightlist
\item
  Обычно используются решающие деревья небольшой глубины (2-6 уровней).
\item
  Основные параметры:

  \begin{itemize}
  \tightlist
  \item
    \texttt{max\_depth} --- максимальная глубина дерева.
  \item
    \texttt{min\_samples\_split} --- минимальное количество объектов для
    разделения узла.
  \end{itemize}
\end{enumerate}

\hypertarget{ux444ux443ux43dux43aux446ux438ux438-ux43fux43eux442ux435ux440ux44c}{%
\subsubsection*{Функции
потерь}\label{ux444ux443ux43dux43aux446ux438ux438-ux43fux43eux442ux435ux440ux44c}}

\begin{enumerate}
\def\labelenumi{\arabic{enumi}.}
\tightlist
\item
  \textbf{Для регрессии:}

  \begin{itemize}
  \tightlist
  \item
    Среднеквадратичная ошибка (MSE):
  \end{itemize}
\end{enumerate}

\[
     L(y, F(x)) = \frac{1}{N} \sum_{i=1}^N (y_i - F(x_i))^2.
\]

\begin{itemize}
\tightlist
\item
  MAE (средняя абсолютная ошибка):
\end{itemize}

\[
     L(y, F(x)) = \frac{1}{N} \sum_{i=1}^N |y_i - F(x_i)|.
\]

\begin{enumerate}
\def\labelenumi{\arabic{enumi}.}
\setcounter{enumi}{1}
\tightlist
\item
  \textbf{Для классификации:}

  \begin{itemize}
  \tightlist
  \item
    Логарифмическая потеря (log-loss):
  \end{itemize}
\end{enumerate}

\[
     L(y, F(x)) = - \frac{1}{N} \sum_{i=1}^N \left( y_i \log p_i + (1 - y_i) \log (1 - p_i) \right).
\]

\hypertarget{ux433ux438ux43fux435ux440ux43fux430ux440ux430ux43cux435ux442ux440ux44b-ux431ux443ux441ux442ux438ux43dux433ux430}{%
\subsubsection*{Гиперпараметры
бустинга}\label{ux433ux438ux43fux435ux440ux43fux430ux440ux430ux43cux435ux442ux440ux44b-ux431ux443ux441ux442ux438ux43dux433ux430}}

\hypertarget{ux43eux441ux43dux43eux432ux43dux44bux435-ux43fux430ux440ux430ux43cux435ux442ux440ux44b}{%
\subsubsection*{Основные
параметры}\label{ux43eux441ux43dux43eux432ux43dux44bux435-ux43fux430ux440ux430ux43cux435ux442ux440ux44b}}

\begin{enumerate}
\def\labelenumi{\arabic{enumi}.}
\tightlist
\item
  \texttt{n\_estimators} --- число деревьев (итераций).

  \begin{itemize}
  \tightlist
  \item
    Большее значение может улучшить качество, но увеличивает риск
    переобучения.
  \end{itemize}
\item
  \texttt{learning\_rate} --- скорость обучения (\(\eta\)).

  \begin{itemize}
  \tightlist
  \item
    Типичные значения: \(0.01 \leq \eta \leq 0.1\).
  \end{itemize}
\item
  \texttt{max\_depth} --- глубина деревьев.
\item
  \texttt{subsample} --- доля объектов для обучения на каждой итерации
  (для стохастического бустинга).
\end{enumerate}

\hypertarget{ux43fux440ux435ux438ux43cux443ux449ux435ux441ux442ux432ux430-ux431ux443ux441ux442ux438ux43dux433ux430}{%
\subsubsection*{Преимущества
бустинга}\label{ux43fux440ux435ux438ux43cux443ux449ux435ux441ux442ux432ux430-ux431ux443ux441ux442ux438ux43dux433ux430}}

\begin{enumerate}
\def\labelenumi{\arabic{enumi}.}
\tightlist
\item
  Высокая точность на большинстве задач.
\item
  Гибкость и возможность адаптации под конкретные задачи.
\item
  Устойчивость к переобучению при правильной настройке гиперпараметров.
\end{enumerate}

\hypertarget{ux43dux435ux434ux43eux441ux442ux430ux442ux43aux438-ux431ux443ux441ux442ux438ux43dux433ux430}{%
\subsubsection*{Недостатки
бустинга}\label{ux43dux435ux434ux43eux441ux442ux430ux442ux43aux438-ux431ux443ux441ux442ux438ux43dux433ux430}}

\begin{enumerate}
\def\labelenumi{\arabic{enumi}.}
\tightlist
\item
  Высокая вычислительная сложность.
\item
  Сложность настройки гиперпараметров.
\item
  Чувствительность к шуму в данных.
\end{enumerate}

\hypertarget{ux441ux43eux432ux440ux435ux43cux435ux43dux43dux44bux435-ux440ux435ux430ux43bux438ux437ux430ux446ux438ux438-ux431ux443ux441ux442ux438ux43dux433ux430}{%
\subsubsection*{Современные реализации
бустинга}\label{ux441ux43eux432ux440ux435ux43cux435ux43dux43dux44bux435-ux440ux435ux430ux43bux438ux437ux430ux446ux438ux438-ux431ux443ux441ux442ux438ux43dux433ux430}}

\hypertarget{xgboost}{%
\subsubsection*{XGBoost}\label{xgboost}}

\begin{enumerate}
\def\labelenumi{\arabic{enumi}.}
\tightlist
\item
  Оптимизированная реализация градиентного бустинга.
\item
  Использует:

  \begin{itemize}
  \tightlist
  \item
    Параллельные вычисления.
  \item
    Регуляризацию для борьбы с переобучением.
  \end{itemize}
\end{enumerate}

\hypertarget{lightgbm}{%
\subsubsection*{LightGBM}\label{lightgbm}}

\begin{enumerate}
\def\labelenumi{\arabic{enumi}.}
\tightlist
\item
  Улучшенная версия бустинга с оптимизацией времени работы.
\item
  Использует метод ``управляемого роста дерева'' (Leaf-wise).
\end{enumerate}

\hypertarget{catboost}{%
\subsubsection*{CatBoost}\label{catboost}}

\begin{enumerate}
\def\labelenumi{\arabic{enumi}.}
\tightlist
\item
  Оптимизирован для категориальных признаков.
\item
  Предотвращает утечки данных в категориальных признаках за счет
  перекрёстного кодирования.
\end{enumerate}

\hypertarget{ux440ux435ux43aux43eux43cux435ux43dux434ux430ux442ux435ux43bux44cux43dux44bux435-ux441ux438ux441ux442ux435ux43cux44b-ux438-ux440ux430ux43dux436ux438ux440ux43eux432ux430ux43dux438ux435}{%
\subsubsection*{Рекомендательные системы и
ранжирование}\label{ux440ux435ux43aux43eux43cux435ux43dux434ux430ux442ux435ux43bux44cux43dux44bux435-ux441ux438ux441ux442ux435ux43cux44b-ux438-ux440ux430ux43dux436ux438ux440ux43eux432ux430ux43dux438ux435}}

Рекомендательные системы --- это методы машинного обучения,
предназначенные для персонализированного подбора контента, продуктов или
услуг для пользователей. Ранжирование --- это процесс сортировки
объектов в соответствии с их релевантностью или важностью для
конкретного запроса или пользователя.

\hypertarget{ux442ux438ux43fux44b-ux440ux435ux43aux43eux43cux435ux43dux434ux430ux442ux435ux43bux44cux43dux44bux445-ux441ux438ux441ux442ux435ux43c}{%
\subsubsection*{Типы рекомендательных
систем}\label{ux442ux438ux43fux44b-ux440ux435ux43aux43eux43cux435ux43dux434ux430ux442ux435ux43bux44cux43dux44bux445-ux441ux438ux441ux442ux435ux43c}}

\hypertarget{ux43aux43eux43bux43bux430ux431ux43eux440ux430ux442ux438ux432ux43dux430ux44f-ux444ux438ux43bux44cux442ux440ux430ux446ux438ux44f-collaborative-filtering}{%
\subsubsection*{Коллаборативная фильтрация (Collaborative
Filtering)}\label{ux43aux43eux43bux43bux430ux431ux43eux440ux430ux442ux438ux432ux43dux430ux44f-ux444ux438ux43bux44cux442ux440ux430ux446ux438ux44f-collaborative-filtering}}

\begin{enumerate}
\def\labelenumi{\arabic{enumi}.}
\tightlist
\item
  \textbf{Основная идея:}

  \begin{itemize}
  \tightlist
  \item
    Основана на предположении, что пользователи с похожими
    предпочтениями будут предпочитать схожие объекты.
  \end{itemize}
\item
  \textbf{Методы:}

  \begin{itemize}
  \tightlist
  \item
    \textbf{User-based:} рекомендации для пользователя на основе
    поведения похожих пользователей.
  \item
    \textbf{Item-based:} рекомендации на основе схожести объектов.
  \end{itemize}
\item
  \textbf{Алгоритмы:}

  \begin{itemize}
  \tightlist
  \item
    Косинусная близость:
  \end{itemize}
\end{enumerate}

\[
     \text{similarity}(u, v) = \frac{\sum_{i=1}^N r_{ui} \cdot r_{vi}}{\sqrt{\sum_{i=1}^N r_{ui}^2} \cdot \sqrt{\sum_{i=1}^N r_{vi}^2}},
\]

где \( r_{ui} \) --- рейтинг пользователя \( u \) для объекта \( i \). -
Матричная факторизация (SVD):

\[
     R \approx U \cdot V^\top,
\]

где \(R \) --- матрица рейтингов, \( U \) и \( V \) --- латентные матрицы
пользователей и объектов.

\begin{enumerate}
\def\labelenumi{\arabic{enumi}.}
\setcounter{enumi}{3}
\tightlist
\item
  \textbf{Преимущества:}

  \begin{itemize}
  \tightlist
  \item
    Высокая точность при наличии достаточного объема данных.
  \item
    Не требует информации о самих объектах.
  \end{itemize}
\item
  \textbf{Недостатки:}

  \begin{itemize}
  \tightlist
  \item
    Проблема холодного старта для новых пользователей или объектов.
  \item
    Высокая вычислительная сложность для больших данных.
  \end{itemize}
\end{enumerate}

\hypertarget{ux43aux43eux43dux442ux435ux43dux442ux43dux430ux44f-ux444ux438ux43bux44cux442ux440ux430ux446ux438ux44f-content-based-filtering}{%
\subsubsection*{Контентная фильтрация (Content-Based
Filtering)}\label{ux43aux43eux43dux442ux435ux43dux442ux43dux430ux44f-ux444ux438ux43bux44cux442ux440ux430ux446ux438ux44f-content-based-filtering}}

\begin{enumerate}
\def\labelenumi{\arabic{enumi}.}
\tightlist
\item
  \textbf{Основная идея:}

  \begin{itemize}
  \tightlist
  \item
    Рекомендации формируются на основе описания объектов и предпочтений
    пользователя.
  \end{itemize}
\item
  \textbf{Методы:}

  \begin{itemize}
  \tightlist
  \item
    Создание профиля пользователя на основе взаимодействия с объектами.
  \item
    Сравнение профиля пользователя с характеристиками объектов.
  \end{itemize}
\item
  \textbf{Алгоритмы:}

  \begin{itemize}
  \tightlist
  \item
    TF-IDF (для текстовых данных):
  \end{itemize}
\end{enumerate}

\[
     \text{tf-idf}(t, d) = \text{tf}(t, d) \cdot \log\frac{N}{\text{df}(t)},
\]


 где \( \text{tf}(t, d) \) — частота термина \( t \) в документе \( d \), \( \text{df}(t) \) — количество документов, содержащих термин \( t \).


\begin{itemize}
\tightlist
\item
  Косинусная близость для сравнения профилей.
\end{itemize}

\begin{enumerate}
\def\labelenumi{\arabic{enumi}.}
\setcounter{enumi}{3}
\tightlist
\item
  \textbf{Преимущества:}

  \begin{itemize}
  \tightlist
  \item
    Хорошо работает для новых пользователей.
  \item
    Простота и интерпретируемость.
  \end{itemize}
\item
  \textbf{Недостатки:}

  \begin{itemize}
  \tightlist
  \item
    Ограничения по разнообразию рекомендаций (рекомендуются только
    похожие объекты).
  \item
    Не учитываются предпочтения других пользователей.
  \end{itemize}
\end{enumerate}

\hypertarget{ux433ux438ux431ux440ux438ux434ux43dux44bux435-ux441ux438ux441ux442ux435ux43cux44b-hybrid-systems}{%
\subsubsection*{Гибридные системы (Hybrid
Systems)}\label{ux433ux438ux431ux440ux438ux434ux43dux44bux435-ux441ux438ux441ux442ux435ux43cux44b-hybrid-systems}}

\begin{enumerate}
\def\labelenumi{\arabic{enumi}.}
\tightlist
\item
  \textbf{Идея:}

  \begin{itemize}
  \tightlist
  \item
    Объединение коллаборативной и контентной фильтрации для повышения
    точности.
  \end{itemize}
\item
  \textbf{Подходы:}

  \begin{itemize}
  \tightlist
  \item
    Комбинирование результатов двух моделей.
  \item
    Обучение единой модели с учетом контентной и коллаборативной
    информации.
  \end{itemize}
\end{enumerate}

\hypertarget{ux440ux430ux43dux436ux438ux440ux43eux432ux430ux43dux438ux435}{%
\subsubsection*{Ранжирование}\label{ux440ux430ux43dux436ux438ux440ux43eux432ux430ux43dux438ux435}}

\hypertarget{ux43eux441ux43dux43eux432ux43dux430ux44f-ux437ux430ux434ux430ux447ux430}{%
\subsubsection*{Основная
задача}\label{ux43eux441ux43dux43eux432ux43dux430ux44f-ux437ux430ux434ux430ux447ux430}}

Ранжирование --- это упорядочивание объектов в соответствии с их
релевантностью для конкретного пользователя или запроса.

\hypertarget{ux442ux438ux43fux44b-ux437ux430ux434ux430ux447-ux440ux430ux43dux436ux438ux440ux43eux432ux430ux43dux438ux44f}{%
\subsubsection*{Типы задач
ранжирования}\label{ux442ux438ux43fux44b-ux437ux430ux434ux430ux447-ux440ux430ux43dux436ux438ux440ux43eux432ux430ux43dux438ux44f}}

\begin{enumerate}
\def\labelenumi{\arabic{enumi}.}
\tightlist
\item
  \textbf{Pointwise:}

  \begin{itemize}
  \tightlist
  \item
    Каждый объект оценивается отдельно.
  \item
    Пример: прогнозирование рейтинга (регрессия).
  \end{itemize}
\item
  \textbf{Pairwise:}

  \begin{itemize}
  \tightlist
  \item
    Модели обучаются на парах объектов, предсказывая, какой объект более
    релевантен.
  \item
    Пример: алгоритмы, использующие SVM или нейронные сети.
  \end{itemize}
\item
  \textbf{Listwise:}

  \begin{itemize}
  \tightlist
  \item
    Учитывается полный список объектов и их относительные ранги.
  \item
    Пример: оптимизация метрик ранжирования, таких как NDCG.
  \end{itemize}
\end{enumerate}

\hypertarget{ux43cux435ux442ux440ux438ux43aux438-ux43eux446ux435ux43dux43aux438-ux440ux430ux43dux436ux438ux440ux43eux432ux430ux43dux438ux44f}{%
\subsubsection*{Метрики оценки
ранжирования}\label{ux43cux435ux442ux440ux438ux43aux438-ux43eux446ux435ux43dux43aux438-ux440ux430ux43dux436ux438ux440ux43eux432ux430ux43dux438ux44f}}

\begin{enumerate}
\def\labelenumi{\arabic{enumi}.}
\tightlist
\item
  \textbf{DCG (Discounted Cumulative Gain):}
\end{enumerate}

\[
   \text{DCG}_p = \sum_{i=1}^p \frac{2^{\text{rel}_i} - 1}{\log_2(i + 1)},
\]

где \( \text{rel}_i \) --- релевантность объекта на позиции \(i \).

\begin{enumerate}
\def\labelenumi{\arabic{enumi}.}
\setcounter{enumi}{2}
\tightlist
\item
  \textbf{NDCG (Normalized DCG):}
\end{enumerate}

\[
   \text{NDCG}_p = \frac{\text{DCG}_p}{\text{IDCG}_p},
\]

где \( \text{IDCG}_p \) --- идеальная DCG.

\begin{enumerate}
\def\labelenumi{\arabic{enumi}.}
\setcounter{enumi}{3}
\tightlist
\item
  \textbf{Precision@K:}
\end{enumerate}

\[
   \text{Precision@K} = \frac{\text{количество релевантных объектов в топ-K}}{K}.
\]

\begin{enumerate}
\def\labelenumi{\arabic{enumi}.}
\setcounter{enumi}{4}
\tightlist
\item
  \textbf{MAP (Mean Average Precision):}
\end{enumerate}

\[
   \text{MAP} = \frac{1}{N} \sum_{q=1}^N \frac{1}{m_q} \sum_{k=1}^{m_q} \text{Precision@k},
\]

где \( N \) --- количество запросов, \( m_q \) --- количество релевантных
объектов для запроса \( q \).

\hypertarget{ux441ux43eux432ux440ux435ux43cux435ux43dux43dux44bux435-ux43fux43eux434ux445ux43eux434ux44b}{%
\subsubsection*{Современные
подходы}\label{ux441ux43eux432ux440ux435ux43cux435ux43dux43dux44bux435-ux43fux43eux434ux445ux43eux434ux44b}}

\subsubsection*{Алгоритмы ранжирования}
\begin{enumerate}
\def\labelenumi{\arabic{enumi}.}
\tightlist
\item
  \textbf{RankNet:} обучение нейронной сети для предсказания порядка в
  парах объектов.
\item
  \textbf{LambdaMART:} ансамблевый метод на основе градиентного
  бустинга, оптимизирующий NDCG.
\end{enumerate}

\hypertarget{ux43fux440ux43eux431ux43bux435ux43cux44b-ux438-ux432ux44bux437ux43eux432ux44b}{%
\subsubsection*{Проблемы и
вызовы}\label{ux43fux440ux43eux431ux43bux435ux43cux44b-ux438-ux432ux44bux437ux43eux432ux44b}}

\begin{enumerate}
\def\labelenumi{\arabic{enumi}.}
\tightlist
\item
  \textbf{Холодный старт:}

  \begin{itemize}
  \tightlist
  \item
    Отсутствие данных для новых пользователей или объектов.
  \item
    Решение: гибридные методы, использование контентных характеристик.
  \end{itemize}
\item
  \textbf{Эффективность и масштабируемость:}

  \begin{itemize}
  \tightlist
  \item
    Обработка большого объема данных.
  \item
    Решение: распределенные системы (Spark, Hadoop).
  \end{itemize}
\item
  \textbf{Предвзятость и разнообразие:}

  \begin{itemize}
  \tightlist
  \item
    Слишком узкие рекомендации.
  \item
    Решение: введение метрик разнообразия и случайности.
  \end{itemize}
\end{enumerate}

\endinput